% 电磁场课程笔记 - 主文档
% 版本：1.0 完整版
% 作者：浮萍
% GitHub：https://github.com/Duckweed-yhb
% 邮箱：3071974740@qq.com

\documentclass{article}
% ==================== 基础宏包 ====================
\usepackage{ctex}                 % 中文支持
\usepackage{amsmath,amssymb}      % 数学公式
\usepackage{esint}                % 环路积分 \oiint 支持
\usepackage{booktabs}             % 表格美化
\usepackage{array}                % 表格扩展
\usepackage{xcolor}               % 颜色
\usepackage{enumitem}             % 列表
\usepackage{geometry}             % 页面设置
\usepackage{titlesec}             % 章节标题
\usepackage{titletoc}             % 目录样式
\usepackage{empheq}               % 公式盒子增强

% ==================== 彩色公式框（统一全局样式） ====================
\usepackage[most]{tcolorbox}
\tcbuselibrary{skins,breakable}

% 统一：重点公式彩色背景框（护眼浅蓝，无黑框）
\newtcolorbox{keyformulabox}{
    colback=blue!8!white,
    colframe=transparent,
    boxrule=0pt,
    arc=2pt,
    left=3mm,
    right=3mm,
    top=2mm,
    bottom=2mm,
    halign=center,
    capture=minipage,
    nobeforeafter,
    width=\linewidth
}

% 兼容旧命令 \keyformula{...} 自动转为彩色框
\newcommand{\keyformula}[1]{%
    \begin{keyformulabox}%
        \displaystyle #1%
    \end{keyformulabox}%
}

% ==================== 页面设置 ====================
\geometry{a4paper,margin=2.5cm}
\setlength{\parindent}{0pt}         % 取消首行缩进
\linespread{1.25}                   % 舒适行距

% ==================== 列表全局样式 ====================
\setlist[itemize]{
    label=\ensuremath{\boldsymbol{\cdot}},
    leftmargin=1.8em,
    itemsep=0.5em,
    parsep=0.3em,
    font=\normalfont
}

% ==================== 章节标题美化 ====================
\titleformat{\section}
  {\centering\Large\bfseries\color{blue!70!black}}
  {\thesection}{1em}{}
\titleformat{\subsection}
  {\large\bfseries\color{black}}
  {\thesubsection}{1em}{}
\titleformat{\subsubsection}
  {\normalsize\bfseries\color{black!80}}
  {\thesubsubsection}{1em}{}

% ==================== 目录美化 ====================
\renewcommand{\contentsname}{\centering\Huge\bfseries 目录}
\titlecontents{section}[0pt]{\bfseries}{\thecontentslabel\quad}{}{\titlerule*[10pt]{.}\contentspage}
\titlecontents{subsection}[1.2em]{\small}{\thecontentslabel\quad}{}{\titlerule*[10pt]{.}\contentspage}

% ==================== 常用数学符号简化定义 ====================
\newcommand{\vecE}{\vec{E}}
\newcommand{\vecD}{\vec{D}}
\newcommand{\vecP}{\vec{P}}
\newcommand{\vecF}{\vec{F}}
\newcommand{\vece}{\vec{e}}
\newcommand{\vecp}{\vec{p}}
\newcommand{\rmd}{\mathrm{d}}        % 标准微分符号
\newcommand{\veczero}{\vec{0}}       % 零矢量
\newcommand{\eps}{\varepsilon}       % 介电常数
\newcommand{\rvec}{\vec{r}}          % 位置矢量

% ==================== 超链接（必须最后加载） ====================
\usepackage{hyperref}
\hypersetup{
    colorlinks=true,
    linkcolor=blue!70!black,
    anchorcolor=blue!70!black,
    citecolor=black,
    urlcolor=blue,
    bookmarksnumbered=true,
    pdfencoding=unicode,
    pdfauthor={浮萍},
    pdftitle={电磁场课程笔记},
    pdfsubject={电磁场与电磁波},
    pdfkeywords={电磁场,静电场,恒定电场,矢量分析}
}

\begin{document}
% 封面
\begin{titlepage}
    \centering
    \vspace*{2cm}
    
    % 主标题
    {\Huge \bfseries 电磁场作业集 \par}
    \vspace{0.8cm}
    
    % 副标题+课程信息
    {\Large \bfseries 考前复习 \par}
    \vspace{0.4cm}
    {\Large 课程代码：\textbf{EE2003} \quad 课程名称：\textbf{电磁场} \par}
    \vspace{0.6cm}
    
    \rule{\linewidth}{0.2mm}
    \vspace{1cm}
    
    % 麦克斯韦名言
    {\itshape \large 电和磁的实验中最明显的现象是，物体之间的相互作用。\par}
    \vspace{0.4cm}
    {\large —— 麦克斯韦 \par}
    \vspace{1.2cm}
    
    % 作者信息
    {\Large 作者：\textbf{浮萍} \par}
    \vspace{0.3cm}
    {\Large 版本：\textbf{1.0} \par}
    \vspace{0.3cm}
    {\Large 邮箱：\textbf{3071974740@qq.com} \par}
    \vspace{0.3cm}
    {\Large GitHub：\textbf{github.com/Duckweed-yhb} \par}
    \vspace{1cm}
    
    % 版权信息
    {\large \copyright\ 2026 \textbf{浮萍} 版权所有 \par}
    {\large \textbf{未经授权，不得复制、传播或修改} \par}
    \vspace{0.5cm}
    
    \vfill
    {\large 生成日期：\textbf{\today} \par}
\end{titlepage}

\section*{版本说明}
\noindent\textbf{版本：V1.0}
\begin{itemize}
    \item \textbf{配套课程：}课程代码\textbf{EE2003}，课程名称\textbf{电磁场}；
    \item \textbf{参考教材：}\textbf{《工程电磁场导论》}，冯慈璋、马西奎主编，西安交通大学出版社2000版，全书内容严格对标教材目录与期末考点；
    \item \textbf{覆盖范围：}全书六大模块：\textbf{第0章矢量分析、第1章静电场、第2章恒定电场、第3章恒定磁场、第4章时变电磁场、第5章准静态电磁场}，对应教材$\S0\sim\S5$全部课内知识点；
    \item \textbf{内容编排：}
    \begin{itemize}
        \item 0.矢量分析：梯度、散度、旋度、高斯公式、斯托克斯公式、正交坐标系算子；
        \item 1.静电场：电场/电位、高斯定理、边值问题、有限差分法(实验内容)、镜像法、电容、静电能量与力（采用库仑定律+虚位移法求电场力）；
        \item 2.恒定电场：导电媒质电流、局外场强、基本方程、分界面衔接条件、恒定电场与静电场类比、电导；
        \item 3.恒定磁场：磁感应强度、安培环路定律、基本方程、磁矢位、标量磁位、磁场能量与力（洛伦兹公式+虚位移法求磁场力）；
        \item 4.时变电磁场：法拉第电磁感应+全电流定律、麦克斯韦方程组、动态位及其积分解、坡印廷矢量、正弦电磁场、电磁辐射；
        \item 5.准静态场：EQS电准静态/MQS磁准静态场、似稳条件、集肤效应、导体交流内阻抗；
    \end{itemize}
    \item \textbf{现存缺点：全篇仅文字与公式，无配套几何示意图、原理插图。}
\end{itemize}

% 目录
\newpage
\tableofcontents
\newpage

% 章节内容
% -------------------------- 矢量分析笔记（子文件） --------------------------
% 说明：本文件为子文件，由主文档 \input{xxx} 引入，不含 document 环境与重复宏包
% 依赖宏包（请在主文档 preamble 中添加）：
% \usepackage{tcolorbox}
% \tcbuselibrary{skins}
% 定义彩色公式框：浅蓝背景、无边框、内边距适中
\newtcolorbox{formulaBox}{
    colback=blue!8!white,   % 浅蓝底色
    colframe=transparent,   % 隐藏边框
    boxrule=0pt,
    halign=center,          % 公式居中
    boxsep=2pt,
    left=4pt,right=4pt,top=3pt,bottom=3pt
}

% -------------------------- 正文 --------------------------
\section{矢量分析}

\subsection{正交坐标系}
\subsubsection{正交坐标系定义}
正交坐标系是指坐标系中各坐标基矢量相互垂直（点积为0）、且为单位矢量的坐标系，满足正交性和归一性，是矢量分析的基础框架。

\subsubsection{三种常用正交坐标系}
\begin{itemize}
    \item \textbf{直角坐标系（笛卡尔坐标系）}：坐标 \((x, y, z)\)，单位矢量 \(\vec{e}_x, \vec{e}_y, \vec{e}_z\)（常矢量），适用于规则几何形状、线性问题。
    \item \textbf{柱坐标系}：坐标 \((\rho, \phi, z)\)，单位矢量 \(\vec{e}_\rho, \vec{e}_\phi, \vec{e}_z\)（\(\vec{e}_\rho, \vec{e}_\phi\) 为变矢量），适用于圆柱对称问题。
    \item \textbf{球坐标系}：坐标 \((r, \theta, \phi)\)，单位矢量 \(\vec{e}_r, \vec{e}_\theta, \vec{e}_\phi\)（均为变矢量），适用于球对称问题。
\end{itemize}

\subsubsection{坐标变换公式}
\begin{itemize}
    \item 直角坐标 \(\leftrightarrow\) 柱坐标
    \[
    \rho = \sqrt{x^2+y^2},\quad \phi = \arctan\frac{y}{x},\quad z=z
    \]
    \item 直角坐标 \(\leftrightarrow\) 球坐标
    \[
    r = \sqrt{x^2+y^2+z^2},\quad \theta = \arccos\frac{z}{r},\quad \phi = \arctan\frac{y}{x}
    \]
\end{itemize}

\vspace{0.8cm}\hrule\vspace{0.8cm}

\subsection{矢量代数}
\subsubsection{矢量基本表示}
矢量为既有大小又有方向的量，直角坐标系中表示为：
\[
\vec{A} = A_x \vec{e}_x + A_y \vec{e}_y + A_z \vec{e}_z
\]
矢量模长：
\[
|\vec{A}| = \sqrt{A_x^2 + A_y^2 + A_z^2}
\]

\subsubsection{矢量点乘（数量积）}
\begin{itemize}
    \item \textbf{定义}：$\vec{A} \cdot \vec{B} = |\vec{A}||\vec{B}|\cos\theta$，$\theta$ 为两矢量夹角。
    \item \textbf{直角坐标（核心公式）}：
    \begin{formulaBox}
    \[
    \vec{A} \cdot \vec{B} = A_xB_x + A_yB_y + A_zB_z
    \]
    \end{formulaBox}
    \item \textbf{性质}：满足交换律、分配律，运算结果为标量。
    \item \textbf{几何意义}：表征两矢量同向程度，正交矢量点积为 $0$。
\end{itemize}

\subsubsection{矢量叉乘（矢量积）}
\begin{itemize}
    \item \textbf{定义}：$\vec{A} \times \vec{B} = |\vec{A}||\vec{B}|\sin\theta \, \vec{e}_n$，运算结果为矢量。
    \item \textbf{直角坐标（行列式，核心公式）}
    \begin{formulaBox}
    \[
    \vec{A} \times \vec{B} = \begin{vmatrix}
    \vec{e}_x & \vec{e}_y & \vec{e}_z \\
    A_x & A_y & A_z \\
    B_x & B_y & B_z
    \end{vmatrix}
    \]
    \end{formulaBox}
    \item \textbf{性质}：反交换律 $\vec{A} \times \vec{B} = -\vec{B} \times \vec{A}$，不满足结合律。
    \item \textbf{几何意义}：模长等于两矢量张成的平行四边形面积。
\end{itemize}

\vspace{0.8cm}\hrule\vspace{0.8cm}

\subsection{标量场与矢量场}
\begin{itemize}
    \item \textbf{标量场}：空间每一点对应一个标量数值，如电势 $\varphi(\vec{r})$、温度 $T(\vec{r})$。
    \item \textbf{等值面（等势面）}：满足 $\varphi(x,y,z) = C$ 的空间曲面，等值面互不相交。
    \item \textbf{矢量场}：空间每一点对应一个矢量，如电场 $\vec{E}(\vec{r})$、磁场 $\vec{H}(\vec{r})$。
    \item \textbf{矢量线方程（核心公式）}：描述矢量场方向的曲线，满足微分关系
    \begin{formulaBox}
    \[
    \frac{dx}{A_x} = \frac{dy}{A_y} = \frac{dz}{A_z}
    \]
    \end{formulaBox}
\end{itemize}

\vspace{0.8cm}\hrule\vspace{0.8cm}

\subsection{标量场的梯度}
\subsubsection{哈密顿算子}
矢量微分算子，直角坐标系形式：
\begin{formulaBox}
\[
\nabla = \frac{\partial}{\partial x}\vec{e}_x + \frac{\partial}{\partial y}\vec{e}_y + \frac{\partial}{\partial z}\vec{e}_z
\]
\end{formulaBox}

\subsubsection{梯度定义与意义}
\begin{itemize}
    \item \textbf{梯度公式（核心）}
    \begin{formulaBox}
    \[
    \nabla \varphi = \frac{\partial \varphi}{\partial x}\vec{e}_x + \frac{\partial \varphi}{\partial y}\vec{e}_y + \frac{\partial \varphi}{\partial z}\vec{e}_z
    \]
    \end{formulaBox}
    \item \textbf{物理意义}：标量场变化率最大的方向，大小等于最大变化率。
    \item \textbf{几何意义}：梯度方向与等值面垂直。
\end{itemize}

\subsubsection{方向导数}
方向导数描述标量场 $\varphi$ 在空间某一点，沿\textbf{指定方向}的变化率。
\begin{itemize}
    \item $l$：任意给定的\textbf{方向矢量}（大小为该方向的步长）。
    \item $\vec{l}_0$：$l$ 方向上的\textbf{单位矢量}，即 $\vec{l}_0 = \dfrac{\vec{l}}{|\vec{l}|}$，模长 $|\vec{l}_0|=1$。
    \item \textbf{计算公式（核心）}
    \begin{formulaBox}
    \[
    \frac{\partial \varphi}{\partial l} = \nabla \varphi \cdot \vec{l}_0
    \]
    \end{formulaBox}
    \item 最大值为 $|\nabla\varphi|$（沿梯度方向），最小值为 $-|\nabla\varphi|$（沿梯度反方向）。
    \item 沿等值面切线方向，方向导数恒为 $0$。
\end{itemize}

\vspace{0.8cm}\hrule\vspace{0.8cm}

\subsection{矢量场的通量与散度}
\subsubsection{通量}
矢量场穿过曲面的流量称为通量：
\[
\Phi = \iint_{S} \vec{A} \cdot \mathrm{d}\vec{S}
\]
闭合曲面通量（外法向为正，常用）：
\begin{formulaBox}
\[
\Phi = \oiint_{S} \vec{A} \cdot \mathrm{d}\vec{S}
\]
\end{formulaBox}

\subsubsection{散度}
散度表征空间微元内场源的发散强度，定义为单位体积闭合曲面的净通量。
\begin{itemize}
    \item \textbf{定义式}
    \[
    \nabla \cdot \vec{A} = \lim_{\Delta V \to 0} \frac{1}{\Delta V}\oiint_{S} \vec{A} \cdot \mathrm{d}\vec{S}
    \]
    \item \textbf{直角坐标（核心公式）}
    \begin{formulaBox}
    \[
    \nabla \cdot \vec{A} = \frac{\partial A_x}{\partial x} + \frac{\partial A_y}{\partial y} + \frac{\partial A_z}{\partial z}
    \]
    \end{formulaBox}
\end{itemize}

\subsubsection{柱坐标散度（重点公式）}
\begin{formulaBox}
\[
\nabla \cdot \vec{A} = \frac{1}{\rho} \frac{\partial (\rho A_\rho)}{\partial \rho} + \frac{1}{\rho} \frac{\partial A_\phi}{\partial \phi} + \frac{\partial A_z}{\partial z}
\]
\end{formulaBox}

\subsubsection{球坐标散度（重点公式）}
\begin{formulaBox}
\[
\nabla \cdot \vec{A} = \frac{1}{r^2} \frac{\partial (r^2 A_r)}{\partial r} 
+ \frac{1}{r \sin\theta} \frac{\partial (A_\theta \sin\theta)}{\partial \theta} 
+ \frac{1}{r \sin\theta} \frac{\partial A_\phi}{\partial \phi}
\]
\end{formulaBox}

\begin{itemize}
    \item \textbf{物理意义}：描述场的源/汇强度，$\nabla\cdot\vec{A}>0$ 为源，$<0$ 为汇。
\end{itemize}

\subsubsection{散度定理（高斯定理，核心定理）}
\begin{formulaBox}
\[
\oiint_{\partial V} \vec{A} \cdot \mathrm{d}\vec{S} = \iiint_V \nabla \cdot \vec{A} \, \mathrm{d} V
\]
\end{formulaBox}
闭合曲面积分 $\iff$ 体积分。

\vspace{0.8cm}\hrule\vspace{0.8cm}

\subsection{矢量场的环量与旋度}
\subsubsection{环量}
矢量场沿闭合曲线的线积分，反映场的旋转特性：
\begin{formulaBox}
\[
\Gamma = \oint_L \vec{A} \cdot \mathrm{d}\vec{l}
\]
\end{formulaBox}

\subsubsection{旋度}
\begin{itemize}
    \item \textbf{直角坐标（行列式，核心）}
    \begin{formulaBox}
    \[
    \nabla \times \vec{A} = \begin{vmatrix}
    \vec{e}_x & \vec{e}_y & \vec{e}_z \\
    \partial_x & \partial_y & \partial_z \\
    A_x & A_y & A_z
    \end{vmatrix}
    \]
    \end{formulaBox}
\end{itemize}

\begin{itemize}
    \item \textbf{柱坐标（行列式，重点）}
    \begin{formulaBox}
    \[
    \nabla \times \vec{A} = \frac{1}{\rho}
    \begin{vmatrix}
    \vec{e}_\rho & \rho\vec{e}_\phi & \vec{e}_z \\
    \frac{\partial}{\partial \rho} & \frac{\partial}{\partial \phi} & \frac{\partial}{\partial z} \\
    A_\rho & \rho A_\phi & A_z
    \end{vmatrix}
    \]
    \end{formulaBox}

    \item \textbf{球坐标（行列式，重点）}
    \begin{formulaBox}
    \[
    \nabla \times \vec{A} = \frac{1}{r^2\sin\theta}
    \begin{vmatrix}
    \vec{e}_r & r\vec{e}_\theta & r\sin\theta\vec{e}_\phi \\
    \frac{\partial}{\partial r} & \frac{\partial}{\partial \theta} & \frac{\partial}{\partial \phi} \\
    A_r & r A_\theta & r\sin\theta A_\phi
    \end{vmatrix}
    \]
    \end{formulaBox}
\end{itemize}

\subsubsection{斯托克斯定理（核心定理）}
\begin{formulaBox}
\[
\oint_L \vec{A} \cdot \mathrm{d}\vec{l} = \iint_S (\nabla \times \vec{A}) \cdot \mathrm{d}\vec{S}
\]
\end{formulaBox}

\vspace{0.8cm}\hrule\vspace{0.8cm}

\subsection{矢量恒等式与亥姆霍兹定理}
\subsubsection{重要矢量恒等式（必考公式）}
\begin{itemize}
    \item 梯度的旋度恒为零：
    \begin{formulaBox}
    \[
    \nabla \times (\nabla \varphi) = 0
    \]
    \end{formulaBox}
    \item 旋度的散度恒为零：
    \begin{formulaBox}
    \[
    \nabla \cdot (\nabla \times \vec{A}) = 0
    \]
    \end{formulaBox}
\end{itemize}

\subsubsection{亥姆霍兹定理}
\begin{itemize}
    \item 若矢量场的散度、旋度及边界条件给定，且场在无穷远处趋于0，则场唯一确定。
    \item 任意矢量场可唯一分解为：无旋场（梯度场）+ 无散场（旋度场）。
    \item 是电磁场理论与流体力学中场的唯一性与分解的基础。
\end{itemize}
% 静电场核心知识点 - 子文件（仅用于 \input 引入）
% 电磁场课程笔记
% 注意：本文件为子文件，无 document 环境，由主文档统一管理
% 依赖宏包（主文档导言区已配置）：
% \usepackage{tcolorbox}
% \tcbuselibrary{skins}
% 彩色公式框定义（同矢量分析，主文档统一定义，此处不再重复）

% -------------------------- 正文 --------------------------
\section{静电场}

\subsection{本章内容概览}
本章系统讲解静电场基础理论与经典求解方法，核心知识点如下：
\begin{itemize}
    \item \textbf{电场强度}、\textbf{标量电位}的定义与相互关系
    \item \textbf{高斯定理}微分形式与积分形式及工程应用
    \item \textbf{静电场基本方程组}、媒质分界面衔接条件
    \item \textbf{泊松/拉普拉斯方程}、边值问题与唯一性定理
    \item \textbf{镜像法}、\textbf{电轴法}的解题思路与典型模型
    \item \textbf{电容}、多导体系统\textbf{部分电容}的分析与计算
    \item \textbf{静电场能量}、\textbf{能量密度}及电场力求解
\end{itemize}

\vspace{0.8cm}\hrule\vspace{0.8cm}

\subsection{电场强度与电位的基本概念}
\subsubsection{研究思路梳理}
以库仑定律为起点，逐层推导静电场基本物理量，逻辑脉络如下：
\begin{enumerate}
    \item \textbf{库仑定律}：描述两点电荷间的相互作用力
    \item \textbf{电场强度}：定义单位正电荷所受电场力
    \item 点电荷、连续分布电荷的电场强度表达式
    \item 由电场无旋性引入\textbf{标量电位}函数
\end{enumerate}
核心结论：\textbf{真空中静电场为无旋场}，满足 $\nabla \times \vec{E} = 0$，因此可引入标量电位简化求解。

\subsubsection{库仑定律}
真空中两点电荷 $q'$、$q$ 间的相互作用力：
\begin{formulaBox}
\[
\vec{F} = \frac{1}{4\pi\varepsilon_0} \frac{q'q}{R^2} \vec{e}_R
= \frac{1}{4\pi\varepsilon_0} \frac{q'q}{|\vec{r}-\vec{r}'|^3} \big(\vec{r}-\vec{r}'\big)
\]
\end{formulaBox}
式中参数说明：
\begin{itemize}
    \item 真空介电常数：$\varepsilon_0 = \dfrac{10^{-9}}{36\pi} \ \mathrm{F/m}$
    \item $\varepsilon$ 为介质绝对介电常数
    \item 适用前提：线性、各向同性、均匀无限大介质
\end{itemize}

\subsubsection{电场强度的定义}
电场的基本作用是对场内电荷施加电场力，电场强度定义为：
\begin{formulaBox}
\[
\vec{E} = \lim_{q_0 \to 0} \frac{\vec{F}}{q_0}
\]
\end{formulaBox}

\paragraph{点电荷的电场强度}
\begin{formulaBox}
\[
\vec{E} = \frac{1}{4\pi\varepsilon_0} \frac{q}{|\vec{r}-\vec{r}'|^3} \big(\vec{r}-\vec{r}'\big)
\]
\end{formulaBox}
物理意义：表征场内某点单位正电荷所受电场力的大小与方向，单位 $\mathrm{V/m}$。

\paragraph{点电荷系电场（场强叠加原理）}
多个点电荷的总电场为各点电荷电场的矢量叠加：
\begin{formulaBox}
\[
\vec{E} = \frac{1}{4\pi\varepsilon_0} \sum_i \frac{q_i}{|\vec{r}-\vec{r}_i'|^3} \big(\vec{r}-\vec{r}_i'\big)
\]
\end{formulaBox}

\subsubsection{连续分布电荷的电场强度}
实际场景中电荷多为连续分布，场强通过体积分/面积分/线积分求解：

\paragraph{体电荷分布 $\rho(\vec{r}')$，单位：$\mathrm{C/m^3}$}
\begin{formulaBox}
\[
\vec{E}(\vec{r}) = \frac{1}{4\pi\varepsilon_0} \int_{V'} \frac{\rho(\vec{r}')}{|\vec{r}-\vec{r}'|^3} \big(\vec{r}-\vec{r}'\big) \mathrm{d} V'
\]
\end{formulaBox}

\paragraph{面电荷分布 $\sigma(\vec{r}')$，单位：$\mathrm{C/m^2}$}
\begin{formulaBox}
\[
\vec{E}(\vec{r}) = \frac{1}{4\pi\varepsilon_0} \int_{S'} \frac{\sigma(\vec{r}')}{|\vec{r}-\vec{r}'|^3} \big(\vec{r}-\vec{r}'\big) \mathrm{d} S'
\]
\end{formulaBox}

\paragraph{线电荷分布 $\tau(\vec{r}')$，单位：$\mathrm{C/m}$}
\begin{formulaBox}
\[
\vec{E}(\vec{r}) = \frac{1}{4\pi\varepsilon_0} \int_{L'} \frac{\tau(\vec{r}')}{|\vec{r}-\vec{r}'|^3} \big(\vec{r}-\vec{r}'\big) \mathrm{d} l'
\]
\end{formulaBox}

\noindent \textbf{补充说明}：面电荷、线电荷可看作体电荷的极限形式，体电荷场强公式具有通用性。

\subsubsection{电位的引出}
直接积分求解电场强度计算繁琐。利用静电场无旋特性，结合矢量微分关系 $\nabla \dfrac{1}{R} = -\dfrac{\vec{e}_R}{R^2}$，可得：
\[
\vec{E} = -\nabla\left( \frac{1}{4\pi\varepsilon_0} \int_{V'} \frac{\rho}{R} \mathrm{d} V' + C \right)
\]

定义\textbf{标量电位函数}：
\[
\varphi = \frac{1}{4\pi\varepsilon_0} \int_{V'} \frac{\rho}{R} \mathrm{d} V' + C
\]
则场强与电位满足核心关系：
\begin{formulaBox}
\[
\vec{E} = -\nabla\varphi
\]
\end{formulaBox}

\paragraph{不同电荷分布的电位公式}
\begin{itemize}
    \item 体电荷：
    \begin{formulaBox}
    \[
    \varphi = \frac{1}{4\pi\varepsilon_0} \int_{V'} \frac{\rho}{R} \mathrm{d} V' + C
    \]
    \end{formulaBox}
    \item 面电荷：
    \begin{formulaBox}
    \[
    \varphi = \frac{1}{4\pi\varepsilon_0} \int_{S'} \frac{\sigma}{R} \mathrm{d} S' + C
    \]
    \end{formulaBox}
    \item 线电荷：
    \begin{formulaBox}
    \[
    \varphi = \frac{1}{4\pi\varepsilon_0} \int_{L'} \frac{\tau}{R} \mathrm{d} l' + C
    \]
    \end{formulaBox}
    \item 点电荷系：
    \begin{formulaBox}
    \[
    \varphi = \frac{1}{4\pi\varepsilon_0} \sum_{k=1}^n \frac{q_k}{R_k} + C
    \]
    \end{formulaBox}
\end{itemize}

\paragraph{电位参考点}
电位包含不定常数 $C$，需人为选取\textbf{电位参考点}（零电位点）确定电位的唯一数值。

\subsubsection{电位与电场强度的积分关系}
单位正电荷沿路径从 $P$ 移动至 $Q$，电场力做功：
\[
W = \int_{P}^{Q} \vec{E} \cdot \mathrm{d}\vec{l}
\]
将 $\vec{E}=-\nabla\varphi$ 代入，得电位差公式：
\begin{formulaBox}
\[
\int_{P}^{Q} \vec{E} \cdot \mathrm{d}\vec{l} = \varphi_P - \varphi_Q
\]
\end{formulaBox}

若取 $Q$ 为电位参考点（$\varphi_Q=0$），则 $P$ 点电位：
\begin{formulaBox}
\[
\varphi_P = \int_{P}^{Q_{\text{ref}}} \vec{E} \cdot \mathrm{d}\vec{l}
\]
\end{formulaBox}
物理意义：将单位正电荷由 $P$ 点移至参考点，电场力所做的功。

\subsubsection{电力线与等位面}
由 $\vec{E} = -\nabla\varphi$ 可推出几何性质：
\begin{itemize}
    \item 电力线与等位面\textbf{处处正交}
    \item 电力线始终指向\textbf{电位降低}的方向
\end{itemize}

\subsubsection{电偶极子}
电偶极子由一对间距极小的等量异号点电荷构成，是分析电介质极化、电磁辐射与干扰的基础模型。

\paragraph{电偶极子电位}
\begin{formulaBox}
\[
\varphi = \frac{ql\cos\theta}{4\pi\varepsilon_0 r^2}
= \frac{1}{4\pi\varepsilon_0} \frac{\vec{p} \cdot \vec{e}_r}{r^2}
\]
\end{formulaBox}
式中 $\vec{p}=q\vec{l}$ 为\textbf{电偶极矩}，方向由负电荷指向正电荷。

\paragraph{电偶极子电场（球坐标系）}
对电位求负梯度，得到电场强度：
\begin{formulaBox}
\[
\vec{E} = \frac{p}{4\pi\varepsilon_0 r^3} \left( 2\cos\theta \, \vec{e}_r + \sin\theta \, \vec{e}_\theta \right)
\]
\end{formulaBox}

\noindent \textbf{基本特性}：电位大小与 $r^2$ 成反比，电场强度大小与 $r^3$ 成反比。

\subsubsection{静电场无旋性等价表述}
静电场无旋是一组互为等价的物理关系：
\begin{formulaBox}
\[
\nabla \times \vec{E} = 0
\iff \vec{E} = -\nabla\varphi
\iff \int_{P}^{Q} \vec{E} \cdot \mathrm{d}\vec{l} = \varphi_P - \varphi_Q
\iff \oint_L \vec{E} \cdot \mathrm{d}\vec{l} = 0
\]
\end{formulaBox}

\vspace{0.8cm}\hrule\vspace{0.8cm}

\subsection{高斯定律及其应用}
静电场属于\textbf{有源无旋场}，高斯定理是描述静电场“有源”特性的核心方程，分为真空形式与电介质形式。

\subsubsection{真空中高斯定理}
静电场由空间中所有电荷共同激发，包含自由电荷与极化电荷两类。
\paragraph{微分形式（散度方程，描述场点局部规律）}
空间任意一点电场强度的散度，等于该点总电荷体密度除以真空介电常数：
\begin{formulaBox}
\[
\nabla \cdot \vec{E} = \frac{\rho}{\varepsilon_0}
\]
\end{formulaBox}
式中 $\rho$ 为该点\textbf{总电荷体密度}，包含自由电荷与极化电荷。

\paragraph{积分形式（电场通量定理，描述宏观闭合曲面规律）}
穿过任意闭合曲面的电场强度总通量，等于该闭合曲面包围的全部电荷量除以真空介电常数：
\begin{formulaBox}
\[
\oint_{S} \vec{E} \cdot \mathrm{d}\vec{S} = \frac{q}{\varepsilon_0}
\]
\end{formulaBox}
该式常用于求解具有球对称、柱对称、平面对称分布的静电场，可大幅简化计算。

\subsubsection{电介质的极化}
电介质在外电场作用下发生极化，产生极化电荷，核心物理量：
\begin{itemize}
    \item \textbf{极化强度}：
    \begin{formulaBox}
    \[
    \vec{P} = \lim_{\Delta V \to 0} \frac{\sum \vec{p}_i}{\Delta V}
    \]
    \end{formulaBox}
    \item \textbf{极化电荷体密度}：
    \begin{formulaBox}
    \[
    \rho_P = -\nabla \cdot \vec{P}
    \]
    \end{formulaBox}
    \item \textbf{极化电荷面密度}：
    \begin{formulaBox}
    \[
    \sigma_P = \vec{P} \cdot \vec{e}_n
    \]
    \end{formulaBox}
\end{itemize}

\subsubsection{电位移矢量与介质中高斯定理}
定义\textbf{电位移矢量}：
\begin{formulaBox}
\[
\vec{D} = \varepsilon_0 \vec{E} + \vec{P}
\]
\end{formulaBox}

对于\textbf{线性、各向同性介质}，满足本构关系：
\begin{formulaBox}
\[
\vec{D} = \varepsilon \vec{E},\quad \varepsilon = \varepsilon_r \varepsilon_0
\]
\end{formulaBox}

介质中高斯定理：
微分形式：
\begin{formulaBox}
\[
\nabla \cdot \vec{D} = \rho_f
\]
\end{formulaBox}
积分形式：
\begin{formulaBox}
\[
\oint_S \vec{D} \cdot \mathrm{d}\vec{S} = q_f
\]
\end{formulaBox}
其中 $\rho_f$、$q_f$ 分别为\textbf{自由电荷体密度}、闭合面内总自由电荷。

\subsubsection{静电平衡下的导体}
导体达到静电平衡后具备三大特性：
\begin{itemize}
    \item 导体内部电场强度 $\vec{E}=0$
    \item 导体为\textbf{等位体}，导体表面为\textbf{等位面}
    \item 自由电荷仅分布于导体外表面
\end{itemize}

\vspace{0.8cm}\hrule\vspace{0.8cm}

\subsection{静电场基本方程组与分界面衔接条件}
不同媒质分界面处场量不连续，微分场方程失效，需由积分形式推导\textbf{分界面衔接条件}。

\subsubsection{电场强度切向衔接条件}
由静电场环路定理 $\displaystyle \oint_L \vec{E} \cdot \mathrm{d}\vec{l}=0$，得切向分量连续：
\begin{formulaBox}
\[
E_{1t} = E_{2t}
\]
\end{formulaBox}
矢量形式：
\begin{formulaBox}
\[
\vec{e}_n \times \big(\vec{E}_2 - \vec{E}_1\big) = 0
\]
\end{formulaBox}

\subsubsection{电位移矢量法向衔接条件}
由介质中高斯定理，得法向分量关系：
\begin{formulaBox}
\[
D_{2n} - D_{1n} = \sigma_f
\]
\end{formulaBox}
矢量形式：
\begin{formulaBox}
\[
\vec{e}_n \cdot \big(\vec{D}_2 - \vec{D}_1\big) = \sigma_f
\]
\end{formulaBox}
$\sigma_f$ 为分界面上的\textbf{自由电荷面密度}。

\subsubsection{两类典型分界面条件}
\paragraph{1. 理想电介质分界面（$\sigma_f=0$）}
此时 $D_{1n}=D_{2n}$，结合本构关系与切向连续条件，推导出\textbf{电场折射定律}：
\begin{formulaBox}
\[
\frac{\tan\alpha_1}{\tan\alpha_2} = \frac{\varepsilon_1}{\varepsilon_2}
\]
\end{formulaBox}
适用：各向同性线性介质、分界面无自由电荷。

\paragraph{2. 导体与介质分界面}
静电平衡时导体内部 $\vec{E}_1=0,\ \vec{D}_1=0$，因此：
\begin{formulaBox}
\[
\begin{cases}
E_{2t} = 0 \\
D_{2n} = \sigma_f
\end{cases}
\]
\end{formulaBox}
介质侧电场仅有法向分量：
\begin{formulaBox}
\[
E_2 = \frac{\sigma_f}{\varepsilon_2}
\]
\end{formulaBox}

\vspace{0.8cm}\hrule\vspace{0.8cm}

\subsection{静电场边值问题与唯一性定理}
\subsubsection{泊松方程与拉普拉斯方程}
将 $\vec{E}=-\nabla\varphi$ 代入高斯定理，得到电位的控制方程——\textbf{泊松方程}：
\begin{formulaBox}
\[
\nabla^2 \varphi = -\frac{\rho}{\varepsilon}
\]
\end{formulaBox}

无源区域（$\rho=0$）方程退化为\textbf{拉普拉斯方程}：
\begin{formulaBox}
\[
\nabla^2 \varphi = 0
\]
\end{formulaBox}

\subsubsection{三类边界条件}
求解边值问题需给定边界条件，分为三类：
\begin{itemize}
    \item 第一类（狄利克雷）：边界电位已知
    \begin{formulaBox}
    \[
    \varphi\big|_S = \varphi_0
    \]
    \end{formulaBox}
    \item 第二类（诺伊曼）：边界电位法向导数已知
    \begin{formulaBox}
    \[
    \dfrac{\partial \varphi}{\partial n}\bigg|_S = f
    \]
    \end{formulaBox}
    \item 混合边界条件：边界同时包含上述两类条件
\end{itemize}

\subsubsection{唯一性定理}
给定完整边界条件后，泊松方程/拉普拉斯方程的解\textbf{唯一存在}。该定理是镜像法、分离变量法等间接解法的理论依据。

\vspace{0.8cm}\hrule\vspace{0.8cm}

\subsection{镜像法与电轴法的求解思路}
\subsubsection{镜像法基本原理}
基于唯一性定理，引入\textbf{虚拟镜像电荷}等效替代导体/介质分界面上的感应电荷，将复杂边值问题转化为简单电荷体系的场叠加问题。

\subsubsection{经典应用模型}
\begin{itemize}
    \item 点电荷 + 接地无限大导体平面
    \item 点电荷 + 接地导体球面
    \item 线电荷 + 平行导体圆柱（\textbf{电轴法}）
\end{itemize}

\vspace{0.8cm}\hrule\vspace{0.8cm}

\subsection{电容与部分电容的计算}
\subsubsection{电容基本概念}
电容是电场对应的电路参数，电容计算本质为静电场求解，流程如下：
\[
Q \xrightarrow{\oint_S \vec{D} \cdot \mathrm{d}\vec{S}} \vec{E} \xrightarrow{\int_A^B \vec{E} \cdot \mathrm{d}\vec{l}} U \implies C = \frac{Q}{U}
\]

\paragraph{电容定义}
两导体系统的电容为导体电荷量与导体间电位差之比：
\begin{formulaBox}
\[
C = \frac{Q}{U}
\]
\end{formulaBox}
单位：法拉($\mathrm{F}$)，工程常用微法($\mathrm{\mu F}$)、皮法($\mathrm{pF}$)。

\paragraph{电容特性}
电容仅由\textbf{导体形状、尺寸、相对位置、介质分布}决定，与带电量、电位差无关。

\paragraph{常用电容器}
高压电力电容器、电子线路用陶瓷电容、电解电容、钽电容等。

\subsubsection{部分电容}
\paragraph{引入背景}
对于三个及以上导体组成的多导体系统，导体间电压会受到所有导体电荷的相互影响，单一电容无法描述该特性，因此引入\textbf{部分电容}。

\paragraph{静电独立系统}
系统内电场仅由内部带电体产生，与外部电荷无关，满足：
\begin{formulaBox}
\[
q_0 + q_1 + q_2 + \dots + q_n = 0
\]
\end{formulaBox}
电位移线仅在系统内部导体之间起止。

\paragraph{电荷与电压的关系}
$n+1$ 导体静电独立系统，各导体电荷满足：
\begin{formulaBox}
\[
\begin{cases}
q_1 = C_{10}U_{10} + C_{12}U_{12} + \dots + C_{1n}U_{1n} \\
\quad\vdots \\
q_n = C_{n1}U_{n1} + C_{n2}U_{n2} + \dots + C_{n0}U_{n0}
\end{cases}
\]
\end{formulaBox}

\begin{itemize}
    \item \textbf{自有部分电容}：
    \begin{formulaBox}
    \[
    C_{k0} = \left. \frac{q_k}{U_{k0}} \right|_{\text{其余电压为0}}
    \]
    \end{formulaBox}
    \item \textbf{互有部分电容}：
    \begin{formulaBox}
    \[
    C_{ki} = \left. \frac{q_k}{U_{ki}} \right|_{\text{其余电压为0}}
    \]
    \end{formulaBox}
\end{itemize}

\paragraph{部分电容性质}
\begin{enumerate}
    \item 所有部分电容恒为正值
    \item 互有部分电容满足互易性：
    \begin{formulaBox}
    \[
    C_{ij}=C_{ji}
    \]
    \end{formulaBox}
    \item 仅与导体结构、介质参数相关
\end{enumerate}

$n+1$ 导体系统中，部分电容总数量：
\begin{formulaBox}
\[
\frac{n(n+1)}{2}
\]
\end{formulaBox}

\paragraph{工程应用}
电路板导线、引脚与接地层之间存在分布部分电容，会造成回路间电磁耦合与干扰。

\vspace{0.8cm}\hrule\vspace{0.8cm}

\subsection{静电能量与电场力的分析}
\subsubsection{静电场能量来源与表达式}
静电场能量由电场建立过程中外力做功转化而来；线性介质中，电场能量仅由最终场分布决定，与充电过程无关。

对电荷逐步建立过程积分，得到静电场总能量：
\begin{formulaBox}
\[
W_e = \frac{1}{2} \int_{V'} \varphi \rho \,\mathrm{d} V'
\]
\end{formulaBox}

不同电荷分布对应的能量公式：
\begin{itemize}
    \item 体电荷系统：
    \begin{formulaBox}
    \[
    W_e = \frac{1}{2} \int_{V'} \varphi \rho \,\mathrm{d} V'
    \]
    \end{formulaBox}
    \item 面电荷系统：
    \begin{formulaBox}
    \[
    W_e = \frac{1}{2} \int_{S'} \varphi \sigma \,\mathrm{d} S'
    \]
    \end{formulaBox}
    \item 多导体系统：
    \begin{formulaBox}
    \[
    W_e = \frac{1}{2}\sum_{k=1}^n \varphi_k q_k
    \]
    \end{formulaBox}
\end{itemize}

\subsubsection{静电能量的空间分布}
结合矢量恒等式、散度定理与 $\vec{E}=-\nabla\varphi$，将能量公式改写为场量积分形式。积分区域拓展至全空间时，无穷远面积分项趋近于零，最终得：
\begin{formulaBox}
\[
W_e = \frac{1}{2} \int_{V} \vec{D} \cdot \vec{E} \,\mathrm{d} V
\]
\end{formulaBox}

\subsubsection{能量密度}
定义\textbf{静电能量密度}（单位体积电场能量）：
\begin{formulaBox}
\[
w_e = \frac{1}{2} \vec{D} \cdot \vec{E}
\]
\end{formulaBox}
结论：静电能量储存在整个电场空间中，电场是能量的载体。

\subsubsection{电场力计算}
工程与理论中常用\textbf{虚位移法（虚功原理）}求解带电导体所受的电场力。
% 恒定电场核心知识点 - 子文件（仅用于 \input 引入）
% 电磁场课程笔记
% 注意：本文件为子文件，无 document 环境，由主文档统一管理
% 依赖宏包（主文档导言区已配置）：
% \usepackage{tcolorbox}
% \tcbuselibrary{skins}
% 彩色公式框定义（同静电场，主文档统一定义，此处不再重复）

% -------------------------- 正文 --------------------------
\section{恒定电场}

\subsection{本章内容概览}
本章系统讲解导电媒质中恒定电场的基础理论，核心知识点如下：
\begin{itemize}
    \item 恒定电场的形成条件与基本特征
    \item 电流、电流密度（体/面/线）与\textbf{欧姆定律的微分形式}
    \item \textbf{焦耳定律微分形式}与功率损耗计算
    \item \textbf{电流连续性方程}（电荷守恒定律）
    \item 恒定电场基本方程组（微分/积分形式）
    \item 分界面衔接条件（标量/矢量形式）与\textbf{电场折射定律}
    \item 导电媒质中恒定电荷的分布
    \item 恒定电场边值问题与\textbf{静电比拟法}
    \item \textbf{电导与部分电导}的计算方法与典型模型
    \item \textbf{接地电阻}、\textbf{跨步电压}的工程应用
\end{itemize}

\vspace{0.8cm}\hrule\vspace{0.8cm}

\subsection{恒定电场的基本概念与特征}
\subsubsection{研究对象与形成条件}
恒定电场的研究对象为\textbf{导电媒质中的恒定电流电场}。
形成条件：
\begin{enumerate}
    \item 导体回路中存在\textbf{恒定电源}（维持电荷分布恒定）
    \item 导电媒质内形成闭合的电流通路
    \item 电荷分布不随时间变化，场量处于动态平衡状态
\end{enumerate}

\subsubsection{核心特征}
恒定电场最根本的特征是：电荷分布不随时间变化，因此场量也不随时间变化，满足：
\begin{formulaBox}
\[
\frac{\partial}{\partial t}(\bullet) = 0
\]
\end{formulaBox}

\noindent 这一特性直接导致：
- 恒定电场与静电场具有形式上的相似性，可采用\textbf{静电比拟法}求解
- 导电媒质内存在恒定电场，导体不再是等位体（除非电流为零）

\vspace{0.8cm}\hrule\vspace{0.8cm}

\subsection{导电媒质中的电流}
\subsubsection{电流的分类}
根据电流形成的物理机制，可分为两类：
\begin{itemize}
    \item \textbf{传导电流}：在导电媒质中（如导体、电解液等）电荷的定向运动形成的电流。
    \item \textbf{运流电流}：在自由空间（如真空等）中，电荷的运动形成的电流，不满足欧姆定律。
\end{itemize}

\subsubsection{电流强度}
电流强度描述了单位时间内通过某一曲面的总电荷量，定义为：
\begin{formulaBox}
\[
I = \frac{\mathrm{d} q}{\mathrm{d} t}
\]
\end{formulaBox}
单位：安培（$\mathrm{A}$），方向规定为正电荷运动的方向。

\paragraph{物理意义说明}
电流强度是一个具有\textbf{通量概念}的物理量，它只描述了每秒通过某一面积的电荷总量，\textbf{无法反映电荷在导体截面上每一点的流动情况}。为了描述电流的空间分布，需要引入\textbf{电流密度矢量}。

\subsubsection{电流密度矢量}
电流密度矢量是描述电流空间分布的基本物理量，其定义为：
\begin{itemize}
    \item \textbf{方向}：与正电荷运动的方向一致；
    \item \textbf{大小}：
    \begin{enumerate}
        \item 单位时间内穿过与电荷运动方向相垂直的单位面积的净电量；
        \item 单位体积内的电荷，以速度 $\vec{v}$ 作匀速运动时，形成的电流密度矢量；
        \item 垂直于电荷运动方向单位面积上通过的电流。
    \end{enumerate}
\end{itemize}

其表达式为：
\begin{formulaBox}
\[
\vec{J} = \rho \vec{v}
\]
\end{formulaBox}
单位：$\mathrm{A/m^2}$。

\subsubsection{不同电荷分布形成的电流}
根据电荷分布的形式，电流可分为三类：

\paragraph{体电流（体电流密度矢量）}
按体密度 $\rho$ 分布的电荷，以速度 $\vec{v}$ 作匀速运动时，形成\textbf{体电流密度}（电流面密度矢量），表达式为：
\begin{formulaBox}
\[
\vec{J} = \rho \vec{v}
\]
\end{formulaBox}
单位：$\mathrm{A/m^2}$。
穿过任意曲面 $S$ 的电流为：
\begin{formulaBox}
\[
I = \iint_S \vec{J} \cdot \mathrm{d}\vec{S}
\]
\end{formulaBox}

\paragraph{面电流（面电流密度矢量）}
按面密度 $\sigma$ 分布的电荷，以速度 $\vec{v}$ 作匀速运动时，形成\textbf{面电流密度}（电流线密度矢量），表达式为：
\begin{formulaBox}
\[
\vec{K} = \sigma \vec{v}
\]
\end{formulaBox}
单位：$\mathrm{A/m}$。
穿过线元 $L$ 的电流为：
\begin{formulaBox}
\[
I = \int_L \vec{K} \cdot \vec{e}_n \, \mathrm{d} l
\]
\end{formulaBox}
其中 $\vec{e}_n$ 为面电流所在表面的法向单位矢量。

\paragraph{线电流}
电荷沿导线轴线流动时，可视为\textbf{线电流}，表达式为：
\begin{formulaBox}
\[
I = \tau v
\]
\end{formulaBox}
其中 $\tau$ 为单位长度电荷，$v$ 为电荷运动速度。

\vspace{0.8cm}\hrule\vspace{0.8cm}

\subsection{欧姆定律与焦耳定律的微分形式}
\subsubsection{欧姆定律的微分形式}
在线性、各向同性导电媒质中，任一点传导电流密度与电场强度成正比：
\begin{formulaBox}
\[
\vec{J}_c = \gamma \vec{E}
\]
\end{formulaBox}
其中：
\begin{itemize}
    \item $\gamma$ 为媒质的\textbf{电导率}，单位：$\mathrm{S/m}$（西门子/米）；
    \item 电阻率 $\rho_r = 1/\gamma$，单位：$\Omega \cdot \mathrm{m}$。
\end{itemize}

\noindent 该式表明：电流密度的方向与电场强度方向一致，大小与电场强度成正比。可与电介质的物性方程 $\vec{D} = \varepsilon \vec{E}$ 类比理解。

\paragraph{与电路欧姆定律的关系}
电路中的欧姆定律 $U=IR$ 是微分形式的宏观表现。对于一段导体，有：
\begin{formulaBox}
\[
I = GU, \quad U = IR
\]
\end{formulaBox}
其中 $G$ 为\textbf{电导}，是电阻 $R$ 的倒数。

\subsubsection{焦耳定律的微分形式}
描述单位体积内的功率损耗（热功率密度），定义为：
\begin{formulaBox}
\[
p = \lim_{\Delta V \to 0} \frac{\Delta P}{\Delta V} = \vec{J}_c \cdot \vec{E}
\]
\end{formulaBox}
单位：$\mathrm{W/m^3}$。

在线性、各向同性导电媒质中，也可表示为：
\begin{formulaBox}
\[
p = \gamma E^2 = \frac{J_c^2}{\gamma}
\]
\end{formulaBox}
又称\textbf{热功率体密度}，描述导体内任一点单位体积的功率损耗。

\noindent 对整个导体体积分，得到总功率损耗：
\begin{formulaBox}
\[
P = \int_V p \,\mathrm{d}V = \int_V \vec{E} \cdot \vec{J}_c \,\mathrm{d}V = I^2 R
\]
\end{formulaBox}

\vspace{0.8cm}\hrule\vspace{0.8cm}

\subsection{电流连续性方程（电荷守恒定律）}
由\textbf{电荷守恒定律}可推导出电流连续性方程，是恒定电场的核心方程之一。

\subsubsection{电荷守恒定律的一般形式}
电荷守恒定律表明：单位时间流出区域 $V$ 的电荷总量等于 $V$ 内电量的减少率，即：
\begin{formulaBox}
\[
\oint_S \vec{J} \cdot \mathrm{d}\vec{S} = -\frac{\partial q}{\partial t}
\]
\end{formulaBox}

\noindent 利用高斯定理可推导出其微分形式：
\begin{formulaBox}
\[
\nabla \cdot \vec{J} + \frac{\partial \rho}{\partial t} = 0
\]
\end{formulaBox}
该式也称为\textbf{电流连续性方程}。

\subsubsection{恒定条件下的简化形式}
对于\textbf{恒定电场}，电荷分布不随时间变化，$\partial \rho/\partial t = 0$，因此方程退化为：
\begin{formulaBox}
\[
\nabla \cdot \vec{J}_c = 0
\]
\end{formulaBox}
积分形式为：
\begin{formulaBox}
\[
\oint_S \vec{J}_c \cdot \mathrm{d}\vec{S} = 0
\]
\end{formulaBox}
物理意义：恒定电流没有头和尾，电流线总是闭合的，流入闭合面的电流等于流出的电流。

\vspace{0.8cm}\hrule\vspace{0.8cm}

\subsection{电源电动势与局外场强}
电源通过局外力（非静电力）维持导体两端的电位差，其作用可等效为\textbf{局外场强} $\vec{E}_e$。

\subsubsection{本节重点}
\begin{itemize}
    \item 区分\textbf{局外场强}和\textbf{电场强度}；
    \item 静电场的电场强度环路线积分为零，电源外恒定电场的环路线积分恒也等于零。
\end{itemize}

\subsubsection{局外力与局外场强}
- \textbf{局外力}：电源中能将正负电荷分离开来的非电磁力 $f_e$；
- \textbf{局外场强}：作用于单位正电荷上的局外力的等效场强，定义为：
\begin{formulaBox}
\[
\vec{E}_e = \lim_{q \to 0} \frac{\vec{f}_e}{q}
\]
\end{formulaBox}
方向由电源负极指向正极，它不是真正的电场强度，仅在引起电流这一点上与电场强度等效。

\subsubsection{含源导电媒质的电流}
在电源内部，电场包括库仑场强 $\vec{E}_c$ 和局外场强 $\vec{E}_e$，因此通过含源导电媒质的电流为：
\begin{formulaBox}
\[
\vec{J} = \gamma (\vec{E}_e + \vec{E}_c)
\]
\end{formulaBox}
在电源以外区域，只存在库仑场强 $\vec{E}_c$。

\subsubsection{电动势定义}
\textbf{电动势}是衡量电源作功能力的物理量，定义为单位正电荷从电源负极运动到正极时，局外力所做的功：
\begin{formulaBox}
\[
\varepsilon = \int_{B}^{A} \vec{E}_e \cdot \mathrm{d}\vec{l}
\]
\end{formulaBox}

\subsubsection{电源内外的场强与环路积分}
\begin{itemize}
    \item 在电源中：总场强 $\vec{E}_\text{T} = \vec{E}_c + \vec{E}_e$；
    \item 在电源外：总场强 $\vec{E}_\text{T} = \vec{E}_c$。
\end{itemize}

若 $L$ 为不经过电源的闭合曲线，则：
\begin{formulaBox}
\[
\oint_L \vec{E}_\text{T} \cdot \mathrm{d}\vec{l} = \oint_L \vec{E}_c \cdot \mathrm{d}\vec{l} = 0
\]
\end{formulaBox}
即恒定电场的环路线积分恒等于零，恒定电场仍是守恒场。

\vspace{0.8cm}\hrule\vspace{0.8cm}

\subsection{恒定电场的基本方程组}
本节核心内容包括：电荷守恒与电流连续性、电场强度的环路线积分、恒定电场的基本方程、分界面衔接条件、导电媒质中恒定电荷的分布、恒定电场的边值问题。

\subsubsection{电源外区域的基本场方程}
恒定电场（电源外）的\textbf{基本方程组}包括微分形式和积分形式：

\begin{itemize}
    \item 电场无旋性：
    \begin{formulaBox}
    \[
    \nabla \times \vec{E} = 0
    \]
    \end{formulaBox}
    积分形式：
    \begin{formulaBox}
    \[
    \oint_L \vec{E} \cdot \mathrm{d}\vec{l} = 0
    \]
    \end{formulaBox}
    物理意义：恒定电场是\textbf{无旋场}，仍是守恒场。

    \item 电流连续性方程（恒定条件）：
    \begin{formulaBox}
    \[
    \nabla \cdot \vec{J}_c = 0
    \]
    \end{formulaBox}
    积分形式：
    \begin{formulaBox}
    \[
    \oint_S \vec{J}_c \cdot \mathrm{d}\vec{S} = 0
    \]
    \end{formulaBox}
    物理意义：恒定电场是\textbf{无源场}，电流线是无头无尾的闭合曲线。

    \item 本构关系（欧姆定律微分形式）：
    \begin{formulaBox}
    \[
    \vec{J}_c = \gamma \vec{E}
    \]
    \end{formulaBox}
    物理意义：导电媒质的物性方程。

    \item 辅助方程（高斯定理与电位移物性方程）：
    \begin{formulaBox}
    \[
    \nabla \cdot \vec{D} = \rho_f, \quad \oint_S \vec{D} \cdot \mathrm{d}\vec{S} = q
    \]
    \end{formulaBox}
    \begin{formulaBox}
    \[
    \vec{D} = \varepsilon \vec{E}
    \]
    \end{formulaBox}
\end{itemize}

\vspace{0.8cm}\hrule\vspace{0.8cm}

\subsection{媒质分界面的衔接条件}
不同导电媒质分界面处，场量满足特定的\textbf{边界条件（衔接条件）}，是求解恒定电场边值问题的关键。

\subsubsection{电场强度的切向分量衔接条件}
由场量闭合曲线的积分方程可得：
\begin{formulaBox}
\[
E_{2t} - E_{1t} = 0 \quad \text{或} \quad \vec{n}_{12} \times (\vec{E}_2 - \vec{E}_1) = 0
\]
\end{formulaBox}
物理意义：分界面两侧电场强度的\textbf{切线分量连续}。

\subsubsection{电流密度的法向分量衔接条件}
由场量闭合曲面的积分方程可得：
\begin{formulaBox}
\[
J_{2n} - J_{1n} = 0 \quad \text{或} \quad \vec{n}_{12} \cdot (\vec{J}_2 - \vec{J}_1) = 0
\]
\end{formulaBox}
物理意义：分界面两侧电流密度的\textbf{法向分量连续}。

\subsubsection{电位移矢量的衔接条件}
电位移矢量 $\vec{D}$ 的衔接条件在恒定电场中依然成立：
\begin{formulaBox}
\[
D_{2n} - D_{1n} = \sigma_f \quad \text{或} \quad \vec{n} \cdot (\vec{D}_2 - \vec{D}_1) = \sigma_f
\]
\end{formulaBox}
物理意义：分界面两侧电位移的法向分量不连续，其差值等于分界面上的\textbf{自由面电荷密度}。

\subsubsection{电位形式的衔接条件}
由 $E_t$ 连续和 $J_n$ 连续可推导出电位的边界条件：
\begin{formulaBox}
\[
\begin{cases}
\varphi_1 = \varphi_2 \\
\gamma_1 \dfrac{\partial \varphi_1}{\partial n} = \gamma_2 \dfrac{\partial \varphi_2}{\partial n}
\end{cases}
\]
\end{formulaBox}

\subsubsection{电场折射定律}
结合上述条件，以及本构关系 $\vec{J} = \gamma \vec{E}$，可推导出恒定电场的\textbf{折射定律}：
\begin{formulaBox}
\[
\frac{\tan\theta_2}{\tan\theta_1} = \frac{\gamma_2}{\gamma_1}
\]
\end{formulaBox}
其中 $\theta_1$、$\theta_2$ 分别为两侧媒质中电流密度（或电场强度）与分界面法线的夹角。

\noindent 对比静电场折射定律：
\begin{formulaBox}
\[
\frac{\tan\alpha_1}{\tan\alpha_2} = \frac{\varepsilon_1}{\varepsilon_2}
\]
\end{formulaBox}

\subsubsection{两类特殊分界面讨论}
\begin{itemize}
    \item \textbf{良导体与不良导体分界面}：假设 $\gamma_2 \gg \gamma_1$，则 $\theta_1 \approx 0$。结论：
    \begin{enumerate}
        \item 从不良导体一侧观察，进入的电流线近似与分界面垂直；
        \item 在不良导体中放入良导体电极，从不良导体一侧看，可认为电流线垂直进入或流出电极表面，电极表面可作为等位面处理。
    \end{enumerate}

    \item \textbf{导体与理想介质分界面}（如架空线与空气）：理想介质中 $\gamma_2=0$，则 $J_{2n}=J_{1n}=0$。结论：
    \begin{enumerate}
        \item 分界面导体侧的电流一定与导体表面平行；
        \item 导体与理想介质分界面上必有电荷分布；
        \item 电场切向分量不为零，导体表面的电场不再与导体表面垂直；一般情况下，法向分量远大于切向分量，可认为 $\vec{E}_2$ 近似垂直于导体表面。
    \end{enumerate}
\end{itemize}

\vspace{0.8cm}\hrule\vspace{0.8cm}

\subsection{导电媒质中恒定电荷的分布}
恒定电场中，高斯定理 $\nabla \cdot \vec{D} = \rho$ 依然成立，结合本构关系可分析导电媒质中的电荷分布。

\subsubsection{均匀导电媒质中的电荷分布}
在各向同性、线性、均匀导电媒质中（$\gamma = \text{常数}$）：
\[
\nabla \cdot \vec{J} = 0 \implies \nabla \cdot (\gamma \vec{E}) = 0 \implies \nabla \cdot \left( \gamma \frac{\vec{D}}{\varepsilon} \right) = 0
\]
由于 $\gamma, \varepsilon$ 为常数，可推出：
\begin{formulaBox}
\[
\rho = 0
\]
\end{formulaBox}
结论：\textbf{均匀导电媒质中无体电荷分布}，电荷只能分布在导电媒质的分界面上。

\subsubsection{不均匀导电媒质中的电荷分布}
在各向同性、线性、不均匀导电媒质中（$\gamma \neq \text{常数}$）：
\[
\rho = \nabla \cdot \vec{D} = \nabla \cdot (\varepsilon \vec{E}) = \nabla \cdot \left( \varepsilon \frac{\vec{J}}{\gamma} \right) = \vec{J} \cdot \nabla \left( \frac{\varepsilon}{\gamma} \right)
\]
由于恒定电场中 $\nabla \cdot \vec{J} = 0$，上式化简为：
\begin{formulaBox}
\[
\rho = \vec{J} \cdot \nabla \left( \frac{\varepsilon}{\gamma} \right)
\]
\end{formulaBox}
结论：不均匀导电媒质中不仅分界面上有面电荷分布，导电媒质内部也有体电荷分布。

\subsubsection{分界面上的自由面电荷密度}
分界面上的自由面电荷密度可由电位移矢量的衔接条件求得：
\begin{formulaBox}
\[
\sigma = D_{2n} - D_{1n} = \varepsilon_2 E_{2n} - \varepsilon_1 E_{1n} = \left( \frac{\varepsilon_2}{\gamma_2} - \frac{\varepsilon_1}{\gamma_1} \right) J_{1n}
\]
\end{formulaBox}

\vspace{0.8cm}\hrule\vspace{0.8cm}

\subsection{恒定电场的边值问题与静电比拟法}
\subsubsection{电位的控制方程}
在恒定电场中，$\nabla \times \vec{E} = 0$，因此可引入标量电位函数 $\varphi$，满足：
\begin{formulaBox}
\[
\vec{E} = -\nabla\varphi
\]
\end{formulaBox}
在各向同性的线性均匀介质中，将 $\vec{J} = \gamma \vec{E} = -\gamma \nabla\varphi$ 代入 $\nabla \cdot \vec{J} = 0$，可得：
\begin{formulaBox}
\[
\nabla^2 \varphi = 0
\]
\end{formulaBox}
即恒定电场的电位满足\textbf{拉普拉斯方程}。

\subsubsection{电位的衔接条件}
分界面上的电位满足：
\begin{formulaBox}
\[
\begin{cases}
\varphi_1 = \varphi_2 \\
\gamma_1 \dfrac{\partial \varphi_1}{\partial n} = \gamma_2 \dfrac{\partial \varphi_2}{\partial n}
\end{cases}
\]
\end{formulaBox}

\subsubsection{静电比拟法的原理}
均匀导电媒质中的恒定电场与均匀介质中的静电场，其电位均满足拉普拉斯方程，且边界条件形式相似，因此可以通过\textbf{静电场的解}来类比求解恒定电场问题。

\subsubsection{对应关系}
\begin{center}
\begin{tabular}{|c|c|}
\hline
\textbf{静电场} & \textbf{恒定电场} \\
\hline
电位 $\varphi$ & 电位 $\varphi$ \\
\hline
电位移矢量 $\vec{D}$ & 电流密度矢量 $\vec{J}$ \\
\hline
电场强度 $\vec{E}$ & 电场强度 $\vec{E}$ \\
\hline
介电常数 $\varepsilon$ & 电导率 $\gamma$ \\
\hline
电荷 $q$ & 电流 $I$ \\
\hline
电容 $C = q/U$ & 电导 $G = I/U$ \\
\hline
\end{tabular}
\end{center}

\vspace{0.8cm}\hrule\vspace{0.8cm}

\subsection{电导与部分电导的计算}
\subsubsection{电导与电阻的定义}
\textbf{电导} $G$ 描述导体的导电能力，定义为：
\begin{formulaBox}
\[
G = \frac{I}{U} = \frac{\int_S \vec{J} \cdot \mathrm{d}\vec{S}}{\int_A^B \vec{E} \cdot \mathrm{d}\vec{l}} = \frac{\int_S \gamma \vec{E} \cdot \mathrm{d}\vec{S}}{\int_A^B \vec{E} \cdot \mathrm{d}\vec{l}}
\]
\end{formulaBox}
单位：西门子（$\mathrm{S}$）。\textbf{电阻} $R$ 是电导的倒数：$R = 1/G$，单位：欧姆（$\Omega$）。

\subsubsection{电导的计算思路}
\begin{itemize}
    \item 已知电流 $I$：$I \to \vec{J} = \gamma \vec{E} \to \vec{E} \to U = \int \vec{E} \cdot \mathrm{d}\vec{l} \to G = I/U$；
    \item 已知电压 $U$：$U \to \vec{E} \to \vec{J} = \gamma \vec{E} \to I = \int \vec{J} \cdot \mathrm{d}\vec{S} \to G = I/U$（不常用）。
\end{itemize}

\subsubsection{绝缘电阻}
\textbf{绝缘电阻}是指用绝缘材料隔开的两个导体之间的直流电阻。绝缘材料中存在绝缘缺陷（如内部气隙、局部开裂、受潮等），在两个导体之间接上高压直流电压时，绝缘材料中会有微弱的泄漏电流流过，其绝缘电阻定义为：
\begin{formulaBox}
\[
R = \frac{\text{两导体间的直流电压 } U}{\text{两导体间的泄漏电流 } I}
\]
\end{formulaBox}
单位：欧姆（$\Omega$），可用于判断绝缘状态。

\subsubsection{静电比拟法计算电导}
利用恒定电场与静电场的相似性，可通过电容 $C$ 类比计算电导 $G$：
\[
C = \frac{Q}{U} = \frac{\int_S \vec{D} \cdot \mathrm{d}\vec{S}}{\int_A^B \vec{E} \cdot \mathrm{d}\vec{l}} = \frac{\varepsilon \int_S \vec{E} \cdot \mathrm{d}\vec{S}}{\int_A^B \vec{E} \cdot \mathrm{d}\vec{l}}
\]
\[
G = \frac{I}{U} = \frac{\int_S \vec{J} \cdot \mathrm{d}\vec{S}}{\int_A^B \vec{E} \cdot \mathrm{d}\vec{l}} = \frac{\gamma \int_S \vec{E} \cdot \mathrm{d}\vec{S}}{\int_A^B \vec{E} \cdot \mathrm{d}\vec{l}}
\]
两者具有相同的积分形式，因此可得：
\begin{formulaBox}
\[
\frac{C}{G} = \frac{\varepsilon}{\gamma}
\]
\end{formulaBox}
即已知电容 $C$ 时，可直接得到电导 $G = \frac{\gamma}{\varepsilon} C$。

\subsubsection{典型电导模型}
常见的电导计算模型包括：
\begin{itemize}
    \item 平板结构电导
    \item 圆柱形导体电导
    \item 球形导体电导
\end{itemize}
计算方法与电容计算类似，可通过静电比拟法求解。

\subsubsection{接地电阻}
\textbf{接地电阻}是电流由接地装置流入大地并扩散至无穷远时所遇到的电阻，包括接地体本身的电阻、接触电阻和大地的散流电阻，常用简化模型为半球形接地体。

\subsubsection{跨步电压}
\textbf{跨步电压}：接地电流流入大地后，地面上相距一步距离（约0.8m）的两点间电位差，是电气安全分析的重要参数。

\vspace{0.8cm}\hrule\vspace{0.8cm}
% 恒定磁场 - 子文件（仅用于 \input 引入）
% 电磁场课程笔记
% 无额外依赖，可直接编译

% -------------------------- 正文 --------------------------
\section{恒定磁场}

\subsection{本章内容概览}
本章系统讲解恒定磁场的基础理论，核心知识点如下：
\begin{itemize}
    \item \textbf{磁感应强度}
    \item \textbf{安培环路定律}
    \item \textbf{恒定磁场基本方程组、分界面上的衔接条件}
    \item \textbf{磁矢位、恒定磁场的边值问题}
    \item \textbf{标量磁位}
    \item \textbf{磁场能量与力}
\end{itemize}

\vspace{0.8cm}\hrule\vspace{0.8cm}

\subsection{磁感应强度}
\subsubsection{引言与研究思路}
本节核心内容为\textbf{真空中恒定磁场的散度}，整体研究思路如下：
\begin{enumerate}
    \item \textbf{安培力定律}
    \item \textbf{磁感应强度}（描述磁场的基本物理量，即\textbf{毕奥-萨伐尔定律}）
    \item 线电流产生的磁感应强度表达式
    \item 体、面电流分布的磁感应强度表达式
    \item 以体电流为例，研究表达式的散度，引出\textbf{磁矢位}
\end{enumerate}

\subsubsection{安培力定律}
安培力定律描述的是\textbf{两个载流线圈之间的作用力}，与静电场中的库仑定律地位相当。

载流线圈 $I_1$ 对 $I_2$ 的作用力为：
\begin{formulaBox}
\[
\vec{F} = \frac{\mu_0}{4\pi} \oint_{L_2} \oint_{L_1} \frac{I_2 d\vec{L}_2 \times \left( I_1 d\vec{L}_1 \times \vec{e}_R \right)}{R^2}
\]
\end{formulaBox}
也可写为：
\[
\vec{F} = \oint_{L_2} I_2 d\vec{L}_2 \times \left( \frac{\mu_0}{4\pi} \oint_{L_1} \frac{I_1 d\vec{L}_1 \times \vec{e}_R}{R^2} \right)
\]

其中：
\begin{itemize}
    \item $\mu_0$ 为\textbf{真空磁导率}，$\mu_0 = 4\pi \times 10^{-7} \, \text{H/m}$；
    \item $R$ 为两个电流元之间的距离；
    \item $\vec{e}_R$ 为从 $I_1 d\vec{L}_1$ 指向 $I_2 d\vec{L}_2$ 的单位矢量。
\end{itemize}

安培力定律与库仑定律相比，公式中存在复杂的方向关系，这种复杂性来源于电流元的矢量性。

\subsubsection{电场与磁场的对比}
磁场与电场的基本属性具有明确的对应关系：
\begin{itemize}
    \item 磁场的基本属性：对运动电荷/电流有力的作用，满足
    \[
    \vec{F} = \oint_{L_2} I_2 d\vec{L}_2 \times \vec{B}
    \]
    \item 电场的基本属性：对电荷有力的作用，满足
    \[
    \vec{F} = q\vec{E}
    \]
\end{itemize}

\subsubsection{毕奥-萨伐尔定律（磁感应强度）}
毕奥-萨伐尔定律给出了电流产生磁场的定量关系，细导线闭合回路激发的磁场为：
\begin{formulaBox}
\[
\vec{B} = \frac{\mu_0}{4\pi} \oint_{L'} \frac{I' d\vec{L}' \times \vec{e}_R}{R^2} = \frac{\mu_0}{4\pi} \oint_{L'} \frac{I' d\vec{L}' \times (\vec{r} - \vec{r}')}{|\vec{r} - \vec{r}'|^3}
\]
\end{formulaBox}

相关说明：
\begin{itemize}
    \item \textbf{磁感应强度}的单位为特斯拉（T），定义为：长1m、电流为1A的直导线在该点处所受到的最大磁场力；
    \item 方向由\textbf{右手螺旋定则}判断：
    \begin{enumerate}
        \item 通电直导线：右手握住直导线，大拇指指向电流方向，弯曲四指方向为磁感应线环绕方向；
        \item 通电螺线管：右手握住螺线管，四指指向电流方向，大拇指指向螺线管内部磁感应线方向；
        \item 环形电流：右手握住环形电流，四指指向电流方向，大拇指指向环形电流内部磁感应线方向。
    \end{enumerate}
\end{itemize}

\subsubsection{不同电流分布的磁感应强度}
根据电流分布形式，毕奥-萨伐尔定律可推广为以下四种形式：
\begin{itemize}
    \item 单个运动点电荷的磁场：
    \[
    \vec{B} = \frac{\mu_0}{4\pi} \frac{q\vec{v} \times \vec{e}_R}{R^2}
    \]
    \item 线电流分布：
    \[
    \vec{B} = \frac{\mu_0}{4\pi} \oint_{L'} \frac{I' d\vec{L}' \times \vec{e}_R}{R^2}
    \]
    \item 面电流分布：
    \[
    \vec{B} = \frac{\mu_0}{4\pi} \int_{S'} \frac{\vec{K} dS' \times \vec{e}_R}{R^2}
    \]
    \item 体电流分布：
    \[
    \vec{B} = \frac{\mu_0}{4\pi} \int_{V'} \frac{\vec{J} dV' \times \vec{e}_R}{R^2}
    \]
\end{itemize}

\subsubsection{矢量磁位的引出}
以体电流的磁场表达式为基础，可通过矢量恒等式引出矢量磁位：
\[
\vec{B} = \frac{\mu_0}{4\pi} \int_{V'} \frac{\vec{J} \times \vec{e}_R}{R^2} dV' = \frac{\mu_0}{4\pi} \int_{V'} \vec{J} \times \left( -\nabla \frac{1}{R} \right) dV'
\]
利用矢量恒等式 $\vec{J} \times \nabla \frac{1}{R} = \nabla \times \left( \frac{\vec{J}}{R} \right)$，可推导得：
\[
\vec{B} = \nabla \times \left( \frac{\mu_0}{4\pi} \int_{V'} \frac{\vec{J}}{R} dV' + \nabla h + C \right)
\]
定义\textbf{矢量磁位} $\vec{A}$：
\begin{formulaBox}
\[
\vec{A} = \frac{\mu_0}{4\pi} \int_{V'} \frac{\vec{J}}{R} dV'
\]
\end{formulaBox}
则磁感应强度与矢量磁位满足：
\[
\vec{B} = \nabla \times \vec{A}
\]

为解决矢量磁位不唯一的问题，引入\textbf{库仑规范}：
\[
\nabla \cdot \vec{A} = 0
\]
并选择无穷远为参考点。

\subsubsection{不同分布电流的矢量磁位}
与磁感应强度对应，不同电流分布的矢量磁位表达式为：
\begin{itemize}
    \item 体电流分布：
    \[
    \vec{A} = \frac{\mu_0}{4\pi} \int_{V'} \frac{\vec{J}}{R} dV'
    \]
    \item 面电流分布：
    \[
    \vec{A} = \frac{\mu_0}{4\pi} \int_{S'} \frac{\vec{K}}{R} dS'
    \]
    \item 线电流分布：
    \[
    \vec{A} = \frac{\mu_0}{4\pi} \oint_{L'} \frac{I}{R} d\vec{L}'
    \]
\end{itemize}

\subsubsection{磁通连续性定理}
由 $\vec{B} = \nabla \times \vec{A}$，利用矢量恒等式 $\nabla \cdot (\nabla \times \vec{A}) = 0$，可得：
\begin{formulaBox}
\[
\nabla \cdot \vec{B} = 0
\]
\end{formulaBox}
积分形式为：
\[
\oint_{S} \vec{B} \cdot d\vec{S} = 0
\]

相关结论：
\begin{itemize}
    \item 恒定磁场是\textbf{无散场}；
    \item 恒定磁场的磁力线是闭合曲线，且与闭合的电流矢量链相链；
    \item 由任一闭合面穿出的净磁通等于零，即穿出的磁通等于穿入的磁通；
    \item 以任意闭合曲线 $L$ 为边界的所有曲面上，具有相同的磁通。
\end{itemize}

\subsubsection{磁感应强度线}
磁感应强度线（$\vec{B}$ 线）的定义：曲线上每一点的切线方向与该点磁感应强度方向一致。

磁感线的微分方程为：
\begin{formulaBox}
\[
\vec{B} \times d\vec{L} = 0 \quad \Rightarrow \quad \frac{B_x}{dx} = \frac{B_y}{dy} = \frac{B_z}{dz}
\]
\end{formulaBox}

\subsubsection{磁偶极子}
磁偶极子是指所围成的面积趋近于零时的载流回路，是研究磁场特性时使用的磁场源模型。

\textbf{磁偶极矩}定义为：
\begin{formulaBox}
\[
\vec{m} = I\vec{S}
\]
\end{formulaBox}
其中 $I$ 为回路电流，$\vec{S}$ 为回路围成的有向面积，方向与电流环绕方向服从右手螺旋关系。

磁偶极子的矢量磁位为：
\[
\vec{A} = \frac{\mu_0}{4\pi} \frac{\vec{m} \times \vec{e}_r}{r^2}
\]
在球坐标系中，其形式为：
\[
\vec{A} = \frac{\mu_0}{4\pi} \frac{IS \sin\theta}{r^2} \vec{e}_\phi
\]
对应的磁感应强度为：
\[
\vec{B} = \nabla \times \vec{A} = \frac{\mu_0 m}{4\pi r^3} \left( 2\cos\theta \vec{e}_r + \sin\theta \vec{e}_\theta \right)
\]

电偶极子电场与磁偶极子磁场的对比：
\begin{itemize}
    \item 电偶极子：
    \[
    \varphi = \frac{1}{4\pi\varepsilon_0} \frac{\vec{p} \cdot \vec{e}_r}{r^2}, \quad \vec{E} = -\nabla\varphi = \frac{p}{4\pi\varepsilon_0 r^3} \left( 2\cos\theta \vec{e}_r + \sin\theta \vec{e}_\theta \right)
    \]
    \item 磁偶极子：
    \[
    \vec{A} = \frac{\mu_0}{4\pi} \frac{\vec{m} \times \vec{e}_r}{r^2}, \quad \vec{B} = \nabla \times \vec{A} = \frac{\mu_0 m}{4\pi r^3} \left( 2\cos\theta \vec{e}_r + \sin\theta \vec{e}_\theta \right)
    \]
\end{itemize}

\vspace{0.8cm}\hrule\vspace{0.8cm}

\subsection{安培环路定律——恒定磁场的旋度方程}
\subsubsection{本节内容概览}
本节核心内容包括：
\begin{enumerate}
    \item \textbf{真空中的安培环路定律}
    \item \textbf{媒质的磁化}
    \item \textbf{磁化强度}
    \item 磁媒质内磁偶极子产生的场
    \item \textbf{磁场强度}
    \item \textbf{媒质中的安培环路定律}
    \item \textbf{磁媒质的物性方程}
\end{enumerate}

\subsubsection{真空中的安培环路定律}
安培环路定律描述了磁场的旋度特性，分为微分形式和积分形式：

\begin{formulaBox}
微分形式：
\[
\nabla \times \vec{B} = \mu_0 \vec{J}_C
\]
积分形式：
\[
\oint_L \vec{B} \cdot d\vec{l} = \mu_0 \iint_S \vec{J}_C \cdot d\vec{S} = \mu_0 \sum I_i
\]
\end{formulaBox}

相关说明：
\begin{itemize}
    \item 磁场是\textbf{有旋场（涡旋场）}，与磁场的散度方程 $\nabla \cdot \vec{B} = 0$ 一致；
    \item 磁感线与电流垂直并相互交链；
    \item 在真空稳恒磁场中，$\vec{B}$ 沿任一闭合路径的积分，等于 $\mu_0$ 乘以该闭合路径所包围面积的各电流的代数和；如果电流的方向和积分回路的绕行方向符合右手螺旋关系，式中的电流取正号；
    \item $\vec{B}$ 对任意闭合曲线 $L$ 的环流不恒等于零，磁场是非保守场，没有磁势的概念；
    \item \textbf{传导电流}：在导电媒质中（如导体、电解液等）电荷的运动形成的电流称为传导电流（自由电流）。
\end{itemize}

\subsubsection{媒质（磁介质）的磁化}
自然界中的所有实体物质都是磁介质。物质由分子或原子组成，分子或原子中任何一个电子都不停地同时参与两种运动，即\textbf{环绕原子核的轨道运动}和\textbf{电子本身的自旋运动}。这两种运动都等效于一个电流分布，因而能产生磁效应。

把分子或原子看成一个整体，分子或原子中各个电子对外界所产生磁效应的总和，可用一个等效的圆电流表示，统称为\textbf{分子电流}。

分子电流具有一定的磁矩，称为\textbf{分子磁矩} $\vec{m}$，其表达式为：
\begin{formulaBox}
\[
\vec{m} = I\vec{S}
\]
\end{formulaBox}
其中：
\begin{itemize}
    \item $I$ 是分子电流强度；
    \item $\vec{S}$ 是分子电流围成的有向曲面的面积，其方向与电流环绕方向服从右手螺旋关系。
\end{itemize}

在没有外磁场作用时，由于热运动，分子磁矩排列是随机的，因此总的磁矩等于零，整块物质对外不显磁性。若把物体放入外磁场中，外磁场将对分子磁矩有转矩作用，其方向与外磁场方向趋于一致，使得分子磁矩的排列比较有序化。此时总磁矩不为零，物质对外呈现磁性，这种现象称为\textbf{媒质（磁介质）的磁化}。

磁媒质中的电流是束缚电流。磁媒质的磁化：磁媒质在外加磁场时，磁偶极矩有转矩作用，其方向与外磁场方向趋于一致，使得磁偶极矩的排列比较有序化的现象称为磁媒质的磁化。

磁化的结果：在磁媒质中出现了宏观的附加电流，称为\textbf{磁化电流}。磁化电流形成附加磁场，从而影响原来磁场的分布。

\subsubsection{磁化强度}
磁化强度是反映媒质被磁化的程度，定义为单位体积内所有磁偶极矩的矢量和：
\begin{formulaBox}
\[
\vec{M} = \lim_{\Delta V \to 0} \frac{\sum_i \vec{m}_i}{\Delta V} \quad (\text{A/m})
\]
\end{formulaBox}
其中 $\vec{m}_i = I\vec{S}$ 是媒质中第 $i$ 个分子的磁偶极矩。

\subsubsection{磁媒质内磁偶极子产生的场}
整个媒质中所有磁偶极矩在场点处产生的矢量磁位为：
\[
\vec{A} = \frac{\mu_0}{4\pi} \int_{V'} \frac{\vec{J}_M(\vec{r}')}{R} dV' + \frac{\mu_0}{4\pi} \oint_{S} \frac{\vec{K}_M(\vec{r}')}{R} dS'
\]
其中：
\begin{itemize}
    \item \textbf{磁化面电流密度}：$\vec{K}_M = \vec{M} \times \vec{e}_n$
    \item \textbf{磁化体电流密度}：$\vec{J}_M = \nabla' \times \vec{M}$
\end{itemize}

磁介质对磁场的影响，可归结为磁化后磁化体电流和磁化面电流在真空所产生的作用。

磁化电流与自由电流来源不同，但从激发磁场这一特性来讲，磁化电流和自由电流没有区别：
\[
\begin{cases}
\nabla \times \vec{B}_C = \mu_0 \vec{J}_C \\
\nabla \cdot \vec{B}_C = 0
\end{cases}
\quad
\begin{cases}
\nabla \times \vec{B}_M = \mu_0 \vec{J}_M \\
\nabla \cdot \vec{B}_M = 0
\end{cases}
\]

\subsubsection{磁场强度}
当有磁介质存在时，磁场可看成由自由电流和磁化电流共同在真空中引起的：
\[
\vec{B} = \vec{B}_C + \vec{B}_M
\]
对其取旋度：
\[
\nabla \times \vec{B} = \nabla \times (\vec{B}_C + \vec{B}_M) = \mu_0 (\vec{J}_C + \vec{J}_M)
\]
代入 $\vec{J}_M = \nabla \times \vec{M}$，整理得：
\[
\nabla \times \left( \frac{\vec{B}}{\mu_0} - \vec{M} \right) = \vec{J}_C
\]

定义\textbf{磁场强度} $\vec{H}$：
\begin{formulaBox}
\[
\vec{H} = \frac{\vec{B}}{\mu_0} - \vec{M} \quad (\text{A/m})
\]
\end{formulaBox}
则上式可写为：
\[
\nabla \times \vec{H} = \vec{J}_C
\]
该式表明，在静磁场中任一点上磁场强度的旋度等于该点的传导电流面密度。

利用斯托克斯定理，可得到积分形式：
\[
\oint_L \vec{H} \cdot d\vec{l} = I
\]
结论：在磁场中，磁场强度沿任一闭合路径的线积分等于穿过该回路所包围面积的自由电流（不包括磁化电流）的代数和。

\subsubsection{磁介质的物性方程}
磁介质的物性方程反映介质的磁化特性：
\begin{formulaBox}
\[
\vec{B} = \mu_0 (\vec{H} + \vec{M})
\]
\end{formulaBox}

在各向同性线性媒质中，磁化强度与磁场强度成正比：
\[
\vec{M} = \chi_M \vec{H}
\]
其中 $\chi_M$ 称为媒质的\textbf{磁化率}，无量纲。代入物性方程可得：
\[
\vec{B} = \mu_0 \vec{H} + \chi_M \mu_0 \vec{H} = \mu_0 (1 + \chi_M) \vec{H} = \mu_0 \mu_r \vec{H} = \mu \vec{H}
\]
其中：
\begin{itemize}
    \item $\mu = \mu_0 \mu_r$ 称为媒质的\textbf{磁导率}，单位 H/m；
    \item $\mu_r = 1 + \chi_M = \mu / \mu_0$ 称为介质的\textbf{相对磁导率}，无量纲。
\end{itemize}

\subsubsection{极化和磁化的对比}
为了方便对比，将介质的极化与磁化特性总结如下：
\begin{center}
\begin{tabular}{|c|c|}
\hline
\textbf{介质的极化} & \textbf{介质的磁化} \\
\hline
极化强度 $\vec{P} = \lim_{\Delta V \to 0} \frac{\sum \vec{p}_i}{\Delta V}$ & 磁化强度 $\vec{M} = \lim_{\Delta V \to 0} \frac{\sum \vec{m}_i}{\Delta V}$ \\
\hline
电偶极矩 $\vec{p}_i = q\vec{d}$ & 磁矩 $\vec{m}_i = I\vec{S}$ \\
\hline
$\sigma_p = \vec{P} \cdot \vec{e}_n, \rho_p = -\nabla \cdot \vec{P}$ & $\vec{K}_m = \vec{M} \times \vec{e}_n, \vec{J}_m = \nabla \times \vec{M}$ \\
\hline
$\nabla \cdot \vec{E} = \frac{\rho_f + \rho_p}{\varepsilon_0}$ & $\nabla \times \vec{B} = \mu_0 (\vec{J}_c + \vec{J}_m)$ \\
\hline
$\vec{D} = \varepsilon_0 \vec{E} + \vec{P}$ & $\vec{B} = \mu_0 (\vec{H} + \vec{M})$ \\
\hline
$\vec{P} = \chi_e \varepsilon_0 \vec{E}$ & $\vec{M} = \chi_m \vec{H}$ \\
\hline
$\vec{D} = \varepsilon \vec{E} = \varepsilon_0 \varepsilon_r \vec{E}$ & $\vec{B} = \mu \vec{H} = \mu_0 \mu_r \vec{H}$ \\
\hline
\end{tabular}
\end{center}

\subsubsection{本节重点}
本节核心重点：
\begin{itemize}
    \item \textbf{真空中和磁媒质中的安培环路定律}（微分、积分形式）——恒定磁场的旋度方程；
    \item 如何应用安培环路定律求磁场。
\end{itemize}

\vspace{0.8cm}\hrule\vspace{0.8cm}

\subsection{恒定磁场基本方程组与分界面衔接条件}
\subsubsection{恒定磁场的基本方程}
恒定磁场的基本方程分为微分形式和积分形式：

\begin{formulaBox}
微分形式：
\[
\begin{cases}
\nabla \times \vec{H} = \vec{J}_C \\
\nabla \cdot \vec{B} = 0
\end{cases}
\]
积分形式：
\[
\begin{cases}
\oint_L \vec{H} \cdot d\vec{l} = I \\
\oint_S \vec{B} \cdot d\vec{S} = 0
\end{cases}
\]
构成方程：
\[
\vec{B} = \mu \vec{H}
\]
\end{formulaBox}

相关说明：
\begin{itemize}
    \item $\nabla \times \vec{H} = \vec{J}_C$：恒定磁场是\textbf{有旋场}，电流密度是磁场的涡旋源；
    \item $\nabla \cdot \vec{B} = 0$：恒定磁场是\textbf{无源场}，磁感应线是无头无尾的闭合曲线，没有磁荷的存在；
    \item $\oint_L \vec{H} \cdot d\vec{l} = I$：恒定磁场的环路线积分等于与积分路径相交链的所有自由电流代数和；
    \item $\oint_S \vec{B} \cdot d\vec{S} = 0$：\textbf{磁通连续性定理}，由任一闭合面穿出的净磁通等于零；
    \item $\vec{B} = \mu \vec{H}$：各向同性线性介质的构成方程。
\end{itemize}

\subsubsection{分界面上的衔接条件}
分界面上的衔接条件描述了磁场在不同介质分界面处的变化规律：

1.  \textbf{磁场强度的衔接条件}：
    \[
    \begin{cases}
    H_{1t} - H_{2t} = K' \\
    \vec{e}_n \times (\vec{H}_2 - \vec{H}_1) = \vec{K}
    \end{cases}
    \]
    说明：由于分界面上电流的存在，分界面两侧磁场强度的切向分量不连续。理论上，在理想导体、超导体表面才可能有\textbf{面电流} $\vec{K}$ 的存在；工程上，唯当电流集肤效应十分显著时，才有可能近似将导体表面的电流分布看作等价的面电流分布。

2.  \textbf{磁感应强度的衔接条件}：
    \[
    \begin{cases}
    B_{2n} - B_{1n} = 0 \\
    \vec{e}_n \cdot (\vec{B}_2 - \vec{B}_1) = 0
    \end{cases}
    \]
    说明：分界面两侧磁感应强度的\textbf{法向分量连续}。

\vspace{0.8cm}\hrule\vspace{0.8cm}

\subsection{磁矢位与恒定磁场的边值问题}
\subsubsection{矢量磁位的微分方程}
根据恒定磁场基本方程的微分形式和物性方程，可推导矢量磁位的微分方程：
\[
\nabla \times \vec{H} = \nabla \times \frac{\vec{B}}{\mu} = \frac{1}{\mu} \nabla \times \nabla \times \vec{A} = \vec{J}
\]
利用矢量恒等式 $\nabla \times \nabla \times \vec{A} = \nabla(\nabla \cdot \vec{A}) - \nabla^2 \vec{A}$，代入库仑规范 $\nabla \cdot \vec{A} = 0$，可得：
\begin{formulaBox}
\[
\nabla^2 \vec{A} = -\mu \vec{J}
\]
\end{formulaBox}
该式表明，矢量磁位满足\textbf{泊松方程}。在没有电流分布的场域中，$\vec{J} = 0$，则方程退化为\textbf{拉普拉斯方程}：
\[
\nabla^2 \vec{A} = 0
\]

\subsubsection{矢量磁位的衔接条件}
根据 $\vec{B} = \nabla \times \vec{A}$，可得分界面上矢量磁位的衔接条件：
\[
\begin{cases}
\vec{A}_1 = \vec{A}_2 \\
\vec{e}_n \times \left( \frac{1}{\mu_2} \nabla \times \vec{A}_2 - \frac{1}{\mu_1} \nabla \times \vec{A}_1 \right) = \vec{K}
\end{cases}
\]

对于平行平面磁场，分界面上的衔接条件可简化为：
\[
\begin{cases}
A_1 = A_2 \\
\frac{1}{\mu_1} \frac{\partial A_1}{\partial n} - \frac{1}{\mu_2} \frac{\partial A_2}{\partial n} = K
\end{cases}
\]

\vspace{0.8cm}\hrule\vspace{0.8cm}

\subsection{标量磁位}
\subsubsection{标量磁位的引出}
如果在所研究的区域内，传导电流处处为零时，$\nabla \times \vec{H} = 0$，此时可以有条件地引入\textbf{标量磁位}：
\begin{formulaBox}
\[
\vec{H} = -\nabla \varphi_m
\]
\end{formulaBox}
其中标量磁位的定义为：
\[
\varphi_{mP} = \int_{P}^{Q} \vec{H} \cdot d\vec{l}
\]
标量磁位是\textbf{多值函数}，其多值性来源于参考点和积分路径的选择。

在均匀媒质中，标量磁位满足拉普拉斯方程：
\[
\nabla^2 \varphi_m = 0
\]

\subsubsection{标量磁位的衔接条件}
\[
\begin{cases}
\varphi_{m1} = \varphi_{m2} \\
\mu_1 \frac{\partial \varphi_{m1}}{\partial n} = \mu_2 \frac{\partial \varphi_{m2}}{\partial n}
\end{cases}
\]

\subsubsection{电位与标量磁位的区别}
\begin{enumerate}
    \item 存在条件：$\varphi_m$ 的提出有条件（传导电流为0）；
    \item 多值性：电位的多值性仅来源于参考点的选择，而磁位的多值性同时受参考点和积分路径的影响；
    \item 物理意义：标量磁位没有直接的物理意义，只是为了方便计算引入的辅助量。
\end{enumerate}

\vspace{0.8cm}\hrule\vspace{0.8cm}

\subsection{磁场能量与力}
\subsubsection{磁场能量与能量密度}
磁场能量来源于磁场建立过程中外力提供的能量。\textbf{磁场能量}的表达式为：
\begin{formulaBox}
\[
W_m = \frac{1}{2} \int_{V'} \vec{A} \cdot \vec{J} dV'
\]
\end{formulaBox}
利用 $\vec{J} = \nabla \times \vec{H}$ 可推导出更通用的形式：
\[
W_m = \frac{1}{2} \int_{V} \vec{B} \cdot \vec{H} dV
\]
该式表明，凡是恒定磁场不为零的空间都储存着磁场能量，场是能量的载体。

\textbf{磁场能量密度}定义为：
\[
w_m = \frac{1}{2} \vec{B} \cdot \vec{H}
\]

\vspace{0.8cm}\hrule\vspace{0.8cm}
% 时变电磁场 - 子文件（仅用于 \input 引入）
% 电磁场课程笔记
% 无额外依赖，可直接编译

% -------------------------- 正文 --------------------------
\section{时变电磁场}

\subsection{本章内容概览}
本章系统讲解时变电磁场的基础理论，核心知识点如下：
\begin{itemize}
    \item \textbf{电磁感应定律和全电流定律}
    \item \textbf{电磁场基本方程组、分界面上的衔接条件}
    \item \textbf{动态位及其积分解}
    \item \textbf{电磁功率流和坡印廷矢量}
    \item \textbf{正弦电磁场}
    \item \textbf{电磁辐射}
\end{itemize}

\vspace{0.8cm}\hrule\vspace{0.8cm}

\subsection{电磁感应定律和全电流定律}
\subsubsection{法拉第定律的推广}
\paragraph{1. 法拉第定律的原始描述}
当与导体回路交链的磁通发生变化时，回路中会产生感应电动势和感应电流，这就是\textbf{法拉第电磁感应定律}。

\begin{formulaBox}
\[
\mathcal{E} = -\frac{\partial \Phi_m}{\partial t} = -\frac{\partial}{\partial t} \int_{S} \vec{B} \cdot d\vec{S}
\]
\end{formulaBox}

负号表示感应电流产生的磁场总是阻碍原磁场的变化。

\paragraph{2. 静止回路的感生电动势}
回路静止不动，磁通变化仅因 $\vec{B}$ 随时间变化时，感应电动势可写为：
\begin{formulaBox}
\[
\mathcal{E} = \oint_{L} \vec{E}_j \cdot d\vec{l} = -\int_{S} \frac{\partial \vec{B}}{\partial t} \cdot d\vec{S}
\]
\end{formulaBox}

该电动势又称为\textbf{感生电动势}，是变压器工作的原理，也称为变压器电势。

\paragraph{3. 麦克斯韦对法拉第定律的推广（感应电场）}
麦克斯韦假说认为：
\begin{itemize}
    \item 变化磁场产生的感应电场 $\vec{E}_j$，不仅存在于导体回路中，而且存在于空间中的任意点上，时变磁场与 $\vec{E}_j$ 相伴存在；
    \item 除了电荷产生电场外，变化的磁场也总要在空间产生电场，由变化磁场产生的电场，称为\textbf{感应电场} $\vec{E}_j$；
    \item 麦克斯韦假说对回路没有任何限制，即不论空间有无导体、有无回路，不论是在真空或媒质中它都适用。
\end{itemize}

法拉第定律的实质：\textbf{时变磁场激发电场（感应电场）}。

\paragraph{4. 库仑场与感应电场的对比}
电荷产生的电场（库仑场）与感应电场的场方程对比：
\begin{center}
\begin{tabular}{|c|c|}
\hline
\textbf{库仑场} & \textbf{感应电场} \\
\hline
$\begin{cases}
\nabla \times \vec{E}_C = 0 \\
\nabla \cdot \vec{E}_C = \dfrac{\rho}{\varepsilon_0}
\end{cases}$
&
$\begin{cases}
\nabla \times \vec{E}_j = -\dfrac{\partial \vec{B}}{\partial t} \\
\nabla \cdot \vec{E}_j = 0
\end{cases}$ \\
\hline
库仑场是\textbf{有源无旋场} & 感应电场是\textbf{有旋无源场（涡旋电场）} \\
\hline
\end{tabular}
\end{center}

自由空间中的总电场 $\vec{E} = \vec{E}_C + \vec{E}_j$ 满足：
\begin{formulaBox}
\[
\begin{cases}
\nabla \times \vec{E} = -\dfrac{\partial \vec{B}}{\partial t} \\
\nabla \cdot \vec{E} = \dfrac{\rho}{\varepsilon_0}
\end{cases}
\]
\end{formulaBox}
该式包含着时变磁场激励电场的重要信息，也被称为\textbf{麦克斯韦第II方程}。

\subsubsection{全电流定律}
\paragraph{1. 问题的提出：安培环路定律的矛盾}
在时变场中，应用安培环路定律会遇到矛盾：
\[
\oint_{L} \vec{H} \cdot d\vec{l} = \int_{S_1} \vec{J} \cdot d\vec{S} = i
\]
\[
\oint_{L} \vec{H} \cdot d\vec{l} = \int_{S_2} \vec{J} \cdot d\vec{S} = 0
\]
这说明传导电流在电容器中是不连续的，安培环路定律在时变场中不再适用。

\paragraph{2. 位移电流的引入}
麦克斯韦假说认为，时变电场也会产生等效的电流——\textbf{位移电流}。
电位移矢量 $\vec{D} = \varepsilon_0 \vec{E} + \vec{P} = \varepsilon \vec{E}$ 的时间变化率可等效为电流密度，称为\textbf{位移电流密度}：
\begin{formulaBox}
\[
\vec{J}_D = \frac{\partial \vec{D}}{\partial t}
\]
\end{formulaBox}

位移电流的物理本质是\textbf{变化的电场}。它解决了时变场中电流不连续的问题，也与电荷守恒定律自洽：
\[
\nabla \cdot \vec{J}_C + \frac{\partial \rho}{\partial t} = \nabla \cdot \vec{J}_C + \frac{\partial}{\partial t} (\nabla \cdot \vec{D}) = \nabla \cdot \left( \vec{J}_C + \frac{\partial \vec{D}}{\partial t} \right) = 0
\]

\paragraph{3. 全电流及其连续性}
全电流是传导电流与位移电流的总和：
\begin{formulaBox}
\[
\vec{J} = \vec{J}_C + \vec{J}_D = \vec{J}_C + \frac{\partial \vec{D}}{\partial t}
\]
\end{formulaBox}

全电流满足连续性：
\[
\nabla \cdot \vec{J} = \nabla \cdot (\vec{J}_C + \vec{J}_D) = 0
\]
即\textbf{全电流是连续的}，且与电荷守恒定律等价。

\paragraph{4. 全电流定律的形式}
引入全电流后，安培环路定律推广为全电流定律：

微分形式（\textbf{麦克斯韦第I方程}）：
\begin{formulaBox}
\[
\nabla \times \vec{H} = \vec{J}_C + \frac{\partial \vec{D}}{\partial t}
\]
\end{formulaBox}

积分形式：
\begin{formulaBox}
\[
\oint_{L} \vec{H} \cdot d\vec{l} = \int_{S} \left( \vec{J}_C + \frac{\partial \vec{D}}{\partial t} \right) \cdot d\vec{S}
\]
\end{formulaBox}

\paragraph{5. 位移电流的磁效应}
麦克斯韦假说认为：位移电流也产生磁场，且与传导电流产生的磁场具有相同的性质。
\[
\begin{cases}
\nabla \times \vec{H}_C = \vec{J}_C \\
\nabla \cdot \vec{B}_C = 0
\end{cases}
\quad
\begin{cases}
\nabla \times \vec{H}_D = \vec{J}_D \\
\nabla \cdot \vec{B}_D = 0
\end{cases}
\]

总磁场的方程为：
\begin{formulaBox}
\[
\begin{cases}
\nabla \times \vec{H} = \vec{J}_C + \dfrac{\partial \vec{D}}{\partial t} \\
\nabla \cdot \vec{B} = 0
\end{cases}
\]
\end{formulaBox}
这是麦克斯韦区别于前人的主要贡献之一，揭示了\textbf{变化的电场也能产生磁场}。

\vspace{0.8cm}\hrule\vspace{0.8cm}

\subsection{电磁场基本方程组、分界面上的衔接条件}
\subsubsection{媒质中的麦克斯韦方程组}
麦克斯韦方程组分为微分形式和积分形式：

\begin{formulaBox}
微分形式：
\[
\begin{cases}
\nabla \times \vec{H} = \vec{J}_C + \dfrac{\partial \vec{D}}{\partial t} \\[6pt]
\nabla \times \vec{E} = -\dfrac{\partial \vec{B}}{\partial t} \\[6pt]
\nabla \cdot \vec{B} = 0 \\[6pt]
\nabla \cdot \vec{D} = \rho_f
\end{cases}
\]
积分形式：
\[
\begin{cases}
\displaystyle \oint_{L} \vec{H} \cdot d\vec{l} = \int_{S} \left( \vec{J}_C + \dfrac{\partial \vec{D}}{\partial t} \right) \cdot d\vec{S} \\[6pt]
\displaystyle \oint_{L} \vec{E} \cdot d\vec{l} = -\int_{S} \dfrac{\partial \vec{B}}{\partial t} \cdot d\vec{S} \\[6pt]
\displaystyle \oint_{S} \vec{B} \cdot d\vec{S} = 0 \\[6pt]
\displaystyle \oint_{S} \vec{D} \cdot d\vec{S} = \int_{V} \rho_f dV
\end{cases}
\]
\end{formulaBox}

同时，宏观电磁场的理论基础还包含\textbf{洛伦兹力公式}：
\begin{formulaBox}
\[
\vec{f} = q\vec{E} + q(\vec{v} \times \vec{B})
\]
\end{formulaBox}

\paragraph{方程组的整体物理意义}
麦克斯韦方程组全面总结了电磁场的运动规律，表明变化的电场和磁场是相互联系、不可分割的统一体。它与洛伦兹力公式一起，构成了宏观电磁场的理论基础。方程组本身是不对称的，原因在于磁场的无源性：不存在磁荷和磁流密度。

\subsubsection{关于麦克斯韦方程组的几点说明}
\begin{enumerate}
    \item \textbf{两个旋度方程是核心方程，第III方程不独立}：
    由法拉第定律可推导磁通连续性方程：
    \[
    \nabla \times \vec{E} = -\frac{\partial \vec{B}}{\partial t}
    \Rightarrow \nabla \cdot (\nabla \times \vec{E}) = \nabla \cdot \left(-\frac{\partial \vec{B}}{\partial t}\right) = 0
    \Rightarrow -\frac{\partial}{\partial t} (\nabla \cdot \vec{B}) = 0
    \Rightarrow \nabla \cdot \vec{B} = C
    \]
    上式在整个时间轴上对任意情况都成立，故 $C=0$，即 $\nabla \cdot \vec{B} = 0$。

    \item \textbf{方程组隐含电荷守恒定律}：
    由安培环路定律和高斯定律可推导出电荷守恒定律：
    \[
    \left.
    \begin{aligned}
    \nabla \times \vec{H} &= \vec{J}_C + \frac{\partial \vec{D}}{\partial t} \\
    \nabla \cdot \vec{D} &= \rho_f
    \end{aligned}
    \right\}
    \Rightarrow \nabla \cdot (\nabla \times \vec{H}) = \nabla \cdot \left( \vec{J}_C + \frac{\partial \vec{D}}{\partial t} \right) = 0
    \Rightarrow \nabla \cdot \vec{J}_C + \frac{\partial \rho_f}{\partial t} = 0
    \]

    \item \textbf{方程组是线性方程组，满足叠加定理条件}。

    \item \textbf{静电场、恒定电场、静磁场是麦克斯韦方程组的特例}：
    令 $\frac{\partial}{\partial t} = 0$，则方程组退化为：
    \begin{formulaBox}
    \[
    \begin{cases}
    \nabla \times \vec{H} = \vec{J}_C \\
    \nabla \times \vec{E} = 0 \\
    \nabla \cdot \vec{B} = 0 \\
    \nabla \cdot \vec{D} = \rho_f
    \end{cases}
    \]
    \end{formulaBox}

    \item \textbf{方程组和各向同性线性媒质的物性方程一起，构成一组完整的电磁学方程}。

    \item \textbf{方程组具有自洽性和完备性}。

    \item \textbf{电磁场的波动性}：
    在无源自由空间中，$\vec{J}_C=0, \rho_f=0$，麦克斯韦方程组为：
    \begin{formulaBox}
    \[
    \begin{cases}
    \nabla \times \vec{H} = \dfrac{\partial \vec{D}}{\partial t} \\
    \nabla \times \vec{E} = -\dfrac{\partial \vec{B}}{\partial t} \\
    \nabla \cdot \vec{B} = 0 \\
    \nabla \cdot \vec{D} = 0
    \end{cases}
    \]
    \end{formulaBox}
    时变电磁场可相互激发，电场和磁场能量相互转换，使电磁场像波浪一样由场源通过任意媒质向远处传播出去，这就是\textbf{电磁波}，是无线电通信的基础。

    在各向同性、线性、均匀媒质中，可推导出\textbf{广义波动方程}：
    \begin{formulaBox}
    \[
    \begin{cases}
    \nabla^2 \vec{H} - \mu\gamma \dfrac{\partial \vec{H}}{\partial t} - \mu\varepsilon \dfrac{\partial^2 \vec{H}}{\partial t^2} = 0 \\[6pt]
    \nabla^2 \vec{E} - \mu\gamma \dfrac{\partial \vec{E}}{\partial t} - \mu\varepsilon \dfrac{\partial^2 \vec{E}}{\partial t^2} = 0
    \end{cases}
    \]
    \end{formulaBox}
    在理想介质中，$\gamma=0$，得到\textbf{齐次波动方程}：
    \begin{formulaBox}
    \[
    \begin{cases}
    \nabla^2 \vec{H} - \dfrac{1}{c^2} \dfrac{\partial^2 \vec{H}}{\partial t^2} = 0 \\[6pt]
    \nabla^2 \vec{E} - \dfrac{1}{c^2} \dfrac{\partial^2 \vec{E}}{\partial t^2} = 0
    \end{cases}
    \quad \text{其中} \quad c = \frac{1}{\sqrt{\mu_0\varepsilon_0}}
    \]
    \end{formulaBox}
    波动是电磁场的基本运动形态。
\end{enumerate}

\subsubsection{时变场的分界面衔接条件}
由麦克斯韦方程组的积分形式可推导出分界面上的边界条件。设介质分界面的法向单位矢量 $\vec{n}_{12}$ 由介质1指向介质2，则衔接条件为：
\begin{formulaBox}
\[
\begin{cases}
\vec{n}_{12} \times (\vec{H}_2 - \vec{H}_1) = \vec{K} \\
\vec{n}_{12} \times (\vec{E}_2 - \vec{E}_1) = 0 \\
\vec{n}_{12} \cdot (\vec{B}_2 - \vec{B}_1) = 0 \\
\vec{n}_{12} \cdot (\vec{D}_2 - \vec{D}_1) = \sigma_f
\end{cases}
\]
\end{formulaBox}

\paragraph{1. 两种理想介质分界面（$\vec{K}=0, \sigma_f=0$）}
此时边界条件简化为：
\[
\begin{cases}
\vec{n}_{12} \times (\vec{H}_2 - \vec{H}_1) = 0 \\
\vec{n}_{12} \times (\vec{E}_2 - \vec{E}_1) = 0 \\
\vec{n}_{12} \cdot (\vec{B}_2 - \vec{B}_1) = 0 \\
\vec{n}_{12} \cdot (\vec{D}_2 - \vec{D}_1) = 0
\end{cases}
\Rightarrow
\begin{cases}
H_{2t} = H_{1t} \\
E_{2t} = E_{1t} \\
B_{2n} = B_{1n} \\
D_{2n} = D_{1n}
\end{cases}
\]
由此可推导出电磁场的\textbf{折射定律}：
\[
\frac{\tan\alpha_1}{\tan\alpha_2} = \frac{\varepsilon_1}{\varepsilon_2}, \quad
\frac{\tan\beta_1}{\tan\beta_2} = \frac{\mu_1}{\mu_2}
\]
其中 $\alpha_1, \alpha_2$ 为电场与分界面法线的夹角，$\beta_1, \beta_2$ 为磁场与分界面法线的夹角。

\paragraph{2. 理想导体表面的边界条件}
设介质1为理想导体，其内部不存在时变电磁场，则边界条件为：
\begin{formulaBox}
\[
\begin{cases}
\vec{n}_{12} \times \vec{H}_2 = \vec{K} \\
\vec{n}_{12} \times \vec{E}_2 = 0 \\
\vec{n}_{12} \cdot \vec{B}_2 = 0 \\
\vec{n}_{12} \cdot \vec{D}_2 = \sigma_f
\end{cases}
\Rightarrow
\begin{cases}
H_{2t} = -K' \\
E_{2t} = 0 \\
B_{2n} = 0 \\
D_{2n} = \sigma_f
\end{cases}
\]
\end{formulaBox}
这表明：在理想导体表面，电场的切向分量为零，磁场的法向分量为零，电力线垂直于导体表面，磁力线平行于导体表面。

\paragraph{3. 衔接条件的物理意义}
连续介质内部的麦克斯韦方程组与媒质分界面上的衔接条件一起，构成了任意区域的电磁场方程。

\vspace{0.8cm}\hrule\vspace{0.8cm}

\subsection{动态位及其积分解}
\subsubsection{动态位的引入与达朗贝尔方程}
在线性、各向同性、均匀介质中，由麦克斯韦方程组引入矢量位 $\vec{A}$ 和标量位 $\varphi$，满足：
\begin{formulaBox}
\[
\vec{B} = \nabla \times \vec{A}, \quad \vec{E} = -\frac{\partial \vec{A}}{\partial t} - \nabla \varphi
\]
\end{formulaBox}

将上述关系代入麦克斯韦第IV方程 $\nabla \cdot \vec{D} = \rho_f$，可得：
\[
\nabla \cdot (\varepsilon \vec{E}) = \rho_f
\Rightarrow \nabla \cdot \left( -\varepsilon \nabla \varphi - \varepsilon \frac{\partial \vec{A}}{\partial t} \right) = \rho_f
\Rightarrow \nabla^2 \varphi + \frac{\partial}{\partial t} (\nabla \cdot \vec{A}) = -\frac{\rho_f}{\varepsilon}
\]

同理，由麦克斯韦第I方程可推导出：
\[
\nabla^2 \vec{A} - \mu\varepsilon \frac{\partial^2 \vec{A}}{\partial t^2} = -\mu \vec{J}_C + \nabla \left( \nabla \cdot \vec{A} + \mu\varepsilon \frac{\partial \varphi}{\partial t} \right)
\]

此时矢量位和标量位耦合在一起，联立求解困难。为了简化，引入\textbf{洛伦兹规范}：
\begin{formulaBox}
\[
\nabla \cdot \vec{A} + \mu\varepsilon \frac{\partial \varphi}{\partial t} = 0
\]
\end{formulaBox}

在此规范下，矢量位和标量位的方程解耦，得到\textbf{达朗贝尔方程}：
\begin{formulaBox}
\[
\begin{cases}
\nabla^2 \vec{A} - \mu\varepsilon \dfrac{\partial^2 \vec{A}}{\partial t^2} = -\mu \vec{J}_C \\[6pt]
\nabla^2 \varphi - \mu\varepsilon \dfrac{\partial^2 \varphi}{\partial t^2} = -\dfrac{\rho_f}{\varepsilon}
\end{cases}
\]
\end{formulaBox}

\subsubsection{达朗贝尔方程的特解（推迟位）}
在无限大均匀媒质中，达朗贝尔方程的特解为：
\begin{formulaBox}
\[
\vec{A}(\vec{r}, t) = \frac{\mu}{4\pi} \int_{V'} \frac{\vec{J}_C(\vec{r}', t - \frac{R}{v})}{R} dV'
\]
\[
\varphi(\vec{r}, t) = \frac{1}{4\pi\varepsilon} \int_{V'} \frac{\rho_f(\vec{r}', t - \frac{R}{v})}{R} dV'
\]
\end{formulaBox}
其中 $R = |\vec{r} - \vec{r}'|$ 为场点与源点的距离，波速 $v = \frac{1}{\sqrt{\mu\varepsilon}}$。

这些解被称为\textbf{推迟位}，其物理意义为：场点 $P$ 在时刻 $t$ 的电磁位，不是由 $t$ 时刻的场源分布决定，而是由 $t' = t - \frac{R}{v}$ 时刻的场源分布决定。相差的时间 $\Delta t = \frac{R}{v}$ 正是电磁作用以速度 $v$ 传播距离 $R$ 所需的时间。

推迟效应说明：
- 电磁扰动是以有限速度传播的，不存在超距作用；
- 空间各点的动态位总是滞后于激励源的变化，电磁波由近及远地向外传播，其波速由媒质特性决定。

\vspace{0.8cm}\hrule\vspace{0.8cm}

\subsection{电磁功率流和坡印廷矢量}
\subsubsection{坡印廷定理}
\paragraph{1. 基本假定与能量密度}
麦克斯韦假定：在均匀各向同性介质中，任意点的电磁能量由电场能量和磁场能量组成，能量密度为：
\begin{formulaBox}
\[
w = w_e + w_m = \frac{1}{2} \vec{E} \cdot \vec{D} + \frac{1}{2} \vec{B} \cdot \vec{H} \quad (\text{单位：J/m}^3)
\]
\end{formulaBox}

体积 $V$ 中电磁场的总能量为：
\[
W = \int_{V} w dV = \int_{V} \left( \frac{1}{2} \vec{E} \cdot \vec{D} + \frac{1}{2} \vec{B} \cdot \vec{H} \right) dV \quad (\text{单位：J})
\]

\paragraph{2. 能量的时间变化率}
能量的时间变化率为：
\[
\frac{\partial W}{\partial t} = \int_{V} \frac{\partial}{\partial t} \left( \frac{1}{2} \vec{E} \cdot \vec{D} + \frac{1}{2} \vec{B} \cdot \vec{H} \right) dV
\]

应用媒质的物性方程 $\vec{D} = \varepsilon \vec{E}$、$\vec{B} = \mu \vec{H}$，可推导出：
\[
\frac{\partial}{\partial t} \left( \frac{1}{2} \vec{D} \cdot \vec{E} \right) = \vec{E} \cdot \frac{\partial \vec{D}}{\partial t}, \quad
\frac{\partial}{\partial t} \left( \frac{1}{2} \vec{B} \cdot \vec{H} \right) = \vec{H} \cdot \frac{\partial \vec{B}}{\partial t}
\]

因此能量变化率可写为：
\[
\frac{\partial W}{\partial t} = \int_{V} \left( \vec{E} \cdot \frac{\partial \vec{D}}{\partial t} + \vec{H} \cdot \frac{\partial \vec{B}}{\partial t} \right) dV
\]

应用麦克斯韦方程组：
\[
\nabla \times \vec{H} = \vec{J}_C + \frac{\partial \vec{D}}{\partial t} \Rightarrow \vec{E} \cdot \frac{\partial \vec{D}}{\partial t} = \vec{E} \cdot (\nabla \times \vec{H}) - \vec{E} \cdot \vec{J}_C
\]
\[
\nabla \times \vec{E} = -\frac{\partial \vec{B}}{\partial t} \Rightarrow \vec{H} \cdot \frac{\partial \vec{B}}{\partial t} = -\vec{H} \cdot (\nabla \times \vec{E})
\]

代入能量变化率表达式，并利用矢量恒等式 $\vec{E} \cdot (\nabla \times \vec{H}) - \vec{H} \cdot (\nabla \times \vec{E}) = \nabla \cdot (\vec{E} \times \vec{H})$，可得：
\[
\frac{\partial W}{\partial t} = \int_{V} \left( \nabla \cdot (\vec{E} \times \vec{H}) - \vec{E} \cdot \vec{J}_C \right) dV
\]

应用高斯散度定理，得到\textbf{坡印廷定理}的最终形式：
\begin{formulaBox}
\[
-\frac{\partial W}{\partial t} = \oint_{A} (\vec{E} \times \vec{H}) \cdot d\vec{A} + \int_{V} \vec{E} \cdot \vec{J}_C dV
\]
\end{formulaBox}

\paragraph{3. 坡印廷定理的物理意义}
坡印廷定理是\textbf{能量守恒定律和转换定律在电磁场中的具体表现}：
- $-\frac{\partial W}{\partial t}$：区域内减少的电磁功率；
- $\oint_{A} \vec{S} \cdot d\vec{A}$：流出区域的电磁功率（能量的流动）；
- $\int_{V} \vec{E} \cdot \vec{J}_C dV$：电磁功率转化为其他形式功率的度量（如焦耳损耗）。

该式表明：电磁场不仅有能量，而且能量是流动的，坡印廷矢量 $\vec{S} = \vec{E} \times \vec{H}$ 表征了这种能量流动的特征。

\subsubsection{坡印廷矢量}
定义\textbf{坡印廷矢量（能流密度矢量）}：
\begin{formulaBox}
\[
\vec{S} = \vec{E} \times \vec{H}
\]
\end{formulaBox}

- 大小：单位为 $\text{W/m}^2$，表示单位时间内通过垂直于传播方向的单位面积的电磁能量；
- 方向：表示电磁能量传播或流动的方向。

其物理意义为：单位时间内，通过垂直于传播方向的单位面积的电磁功率。

\subsubsection{典型实例分析}
\paragraph{1. 长直同轴电缆的能流}
设内导体半径为 $a$，外导体内半径为 $b$，内外导体间电压为 $U$，流过电流为 $I$。

- \textbf{理想导体情况}：
  内导体内电场为0，能量仅在内外导体之间流动。内外导体之间：
  \[
  \vec{E} = \frac{U}{\rho \ln(b/a)} \vec{e}_\rho, \quad
  \vec{H} = \frac{I}{2\pi\rho} \vec{e}_\varphi
  \]
  坡印廷矢量：
  \[
  \vec{S} = \vec{E} \times \vec{H} = \frac{UI}{2\pi\rho^2 \ln(b/a)} \vec{e}_z
  \]
  传输功率：
  \[
  P = \int_{A} \vec{S} \cdot d\vec{A} = UI
  \]
  结论：电磁能量在导体之间沿轴线流动，由电源流向负载，导体的作用是建立场，为能量传输指明方向。

- \textbf{非理想导体情况}：
  内导体内部存在轴向电场 $E_z = \frac{I}{\pi a^2 \gamma}$，坡印廷矢量有径向分量，指向内导体：
  \[
  S_\rho = E_z \times H_\varphi = -\frac{I^2 r}{2\pi a}
  \]
  这部分能量在导体内部转化为焦耳热损耗，单位长度损耗功率为 $I^2 r$。

\paragraph{2. 圆柱形平板电容器的能流}
缓慢充电时，电容器内部的能量增加是从侧面流进去的，而不是通过极板上的导线流入。这进一步说明导线不是能量的载体，能量是由场传输的。

\vspace{0.8cm}\hrule\vspace{0.8cm}

\subsection{正弦电磁场}
\subsubsection{正弦电磁场的定义与研究意义}
以一定频率做正弦变化的场，称为\textbf{正弦电磁场}。

研究正弦电磁场的意义：一般情况下，非正弦变化的时变场可以应用傅里叶级数将它分解成稳态场和频率分量各不相同的正弦电磁场，来分别加以研究。

\subsubsection{正弦电磁场的复数表示法}
在直角坐标系中，随时间作正弦变化的电场强度的一般形式为：
\[
\vec{E}(x,y,z,t) = E_{xm}(x,y,z)\cos(\omega t + \phi_x)\vec{e}_x + E_{ym}(x,y,z)\cos(\omega t + \phi_y)\vec{e}_y + E_{zm}(x,y,z)\cos(\omega t + \phi_z)\vec{e}_z
\]

引入\textbf{电场强度复矢量} $\dot{\vec{E}}(x,y,z)$，其定义为：
\begin{formulaBox}
\[
\vec{E}(x,y,z,t) = \text{Re}\left[ \dot{\vec{E}}(x,y,z) \sqrt{2} e^{j\omega t} \right]
\]
其中：
\[
\dot{\vec{E}}(x,y,z) = \dot{E}_x \vec{e}_x + \dot{E}_y \vec{e}_y + \dot{E}_z \vec{e}_z = \frac{1}{\sqrt{2}} E_{xm} e^{j\phi_x} \vec{e}_x + \frac{1}{\sqrt{2}} E_{ym} e^{j\phi_y} \vec{e}_y + \frac{1}{\sqrt{2}} E_{zm} e^{j\phi_z} \vec{e}_z
\]
\end{formulaBox}

复数法的优势：对时间的求导运算可以转化为乘积运算：
\[
\frac{\partial \vec{E}(x,y,z,t)}{\partial t} = \text{Re}\left[ j\omega \dot{\vec{E}}(x,y,z) \sqrt{2} e^{j\omega t} \right]
\]

\subsubsection{坡印廷定理的复数形式}
对于正弦电磁场，坡印廷定理的复数形式为：
\begin{formulaBox}
\[
-\oint_{A} (\dot{\vec{E}} \times \dot{\vec{H}}^*) \cdot d\vec{A} = \int_{V} \dot{\vec{J}}_C^* \cdot \dot{\vec{E}} dV + j\omega \int_{V} (\dot{\vec{B}} \cdot \dot{\vec{H}}^* - \dot{\vec{E}} \cdot \dot{\vec{D}}^*) dV
\]
\end{formulaBox}

定义\textbf{复坡印廷矢量} $\dot{\vec{S}} = \dot{\vec{E}} \times \dot{\vec{H}}^*$，流入闭合面 $A$ 的复功率为：
\[
\widetilde{S} = -\oint_{A} \dot{\vec{S}} \cdot d\vec{A}
\]

- 有功功率（平均功率）：$P = \text{Re}[\widetilde{S}] = \int_{V} \dot{\vec{J}}_C^* \cdot \dot{\vec{E}} dV$，表示体积 $V$ 内媒质焦耳热损耗的平均值；
- 无功功率：$Q = \text{Im}[\widetilde{S}] = \omega \int_{V} (\dot{\vec{B}} \cdot \dot{\vec{H}}^* - \dot{\vec{E}} \cdot \dot{\vec{D}}^*) dV$，表示体积 $V$ 内电磁能量的交换。

引入复坡印廷矢量的目的是便于直接用矢量计算能流密度的大小和方向。复数形式侧重在有功功率、无功功率或复功率守恒，与时域形式没有一一对应的关系。

\subsubsection{达朗贝尔方程的复数形式及其解}
对于正弦电磁场，达朗贝尔方程的复数形式为：
\begin{formulaBox}
\[
\begin{cases}
\nabla^2 \dot{\vec{A}} + \beta^2 \dot{\vec{A}} = -\mu \dot{\vec{J}}_C \\[6pt]
\nabla^2 \dot{\varphi} + \beta^2 \dot{\varphi} = -\dfrac{\dot{\rho}_f}{\varepsilon}
\end{cases}
\]
其中相位常数（\textbf{波数}）$\beta = \omega\sqrt{\mu\varepsilon}$，单位为 rad/m。
\end{formulaBox}

在无限大均匀媒质中，动态位解的复数形式（滞后位）为：
\begin{formulaBox}
\[
\dot{\vec{A}} = \frac{\mu}{4\pi} \int_{V'} \frac{\dot{\vec{J}}_C e^{-j\beta R}}{R} dV'
\]
\[
\dot{\varphi} = \frac{1}{4\pi\varepsilon} \int_{V'} \frac{\dot{\rho}_f e^{-j\beta R}}{R} dV'
\]
\end{formulaBox}
其中 $e^{-j\beta R}$ 表示场点变化滞后于源点变化的相位差为 $\beta R$，该项的存在表明了场的波动性。

电场、磁场与动态位的关系在复数形式下为：
\[
\begin{cases}
\dot{\vec{B}} = \nabla \times \dot{\vec{A}} \\
\dot{\vec{E}} = -j\omega \dot{\vec{A}} - \nabla \dot{\varphi} = -j\omega \dot{\vec{A}} + \frac{\nabla(\nabla \cdot \dot{\vec{A}})}{j\omega\mu\varepsilon}
\end{cases}
\]
洛伦兹规范的复数形式为：$\nabla \cdot \dot{\vec{A}} + j\omega\mu\varepsilon \dot{\varphi} = 0$。因此，只要求出 $\dot{\vec{A}}$，就可计算出电场和磁场。

\vspace{0.8cm}\hrule\vspace{0.8cm}

\subsection{电磁辐射}
\subsubsection{电磁辐射的基本概念}
电磁能量脱离源而单独存在于空间中，这种现象称为\textbf{电磁辐射}。

\paragraph{电磁辐射的普遍性}
任何物体都有辐射，其本质是内部带电粒子的运动激发电磁场向外传播。日常接触的红外线、微波等都属于电磁辐射，而只有频率极高的电离辐射（如紫外线、X射线、$\gamma$ 射线）才可能对人体健康产生致癌影响。

\subsubsection{单元偶极子天线辐射}
长度 $l$ 远小于波长的线电流元 $I\mathrm{d}\vec{l}$ 是辐射场源的基本组成单元，称为\textbf{电基本振子}，又称单元偶极子。

设电流元为 $\dot{I}\Delta l$，则其矢量位为：
\[
\dot{\vec{A}}(\vec{r}) = \frac{\mu_0 \dot{I} \Delta l}{4\pi r} e^{-j\beta r} \vec{e}_z
\]

代入电场和磁场与动态位的关系，可得单元偶极子在自由空间中激发的电磁场：
\begin{formulaBox}
\[
\begin{cases}
\dot{E}_r = \sqrt{\dfrac{\mu_0}{\varepsilon_0}} \dfrac{\dot{I} \Delta l \beta^2}{2\pi} e^{-j\beta r} \cos\theta \left[ \dfrac{1}{(\beta r)^2} - j\dfrac{1}{(\beta r)^3} \right] \\[6pt]
\dot{E}_\theta = \sqrt{\dfrac{\mu_0}{\varepsilon_0}} \dfrac{\dot{I} \Delta l \beta^2}{4\pi} e^{-j\beta r} \sin\theta \left[ \dfrac{1}{(\beta r)^2} + j\dfrac{1}{\beta r} - j\dfrac{1}{(\beta r)^3} \right] \\[6pt]
\dot{H}_\varphi = \dfrac{\dot{I} \Delta l \beta^2}{4\pi} e^{-j\beta r} \sin\theta \left[ \dfrac{1}{(\beta r)^2} + j\dfrac{1}{\beta r} \right] \\[6pt]
\dot{E}_\varphi = \dot{H}_r = \dot{H}_\theta = 0
\end{cases}
\]
\end{formulaBox}

\subsubsection{辐射场的分区特性}
根据场点与源点的距离 $r$ 与波长 $\lambda$ 的关系，可将场分为近区和远区。

\paragraph{1. 近区电磁场（感应区，$\beta r \ll 1$）}
此时推迟效应可以忽略，场的分布特性与恒定电磁场类似：
\begin{formulaBox}
\[
\begin{cases}
\dot{E}_r = -j\dfrac{\dot{I} \Delta l}{2\pi\omega\varepsilon_0 r^3} \cos\theta = \dfrac{\dot{p}}{2\pi\varepsilon_0 r^3} \cos\theta \\[6pt]
\dot{E}_\theta = -j\dfrac{\dot{I} \Delta l}{4\pi\omega\varepsilon_0 r^3} \sin\theta = \dfrac{\dot{p}}{4\pi\varepsilon_0 r^3} \sin\theta \\[6pt]
\dot{H}_\varphi = \dfrac{\dot{I} \Delta l}{4\pi r^2} \sin\theta
\end{cases}
\]
其中电偶极矩 $\dot{p} = \dfrac{\dot{I} \Delta l}{j\omega}$。
\end{formulaBox}

近区场的特点：
- 电场和磁场相位相差90度，平均能流为零，主要是无功功率；
- 只考虑电能和磁能的交换和振荡，忽略向外辐射的电磁能量，即认为没有能量向外辐射。

\paragraph{2. 远区电磁场（辐射区，$\beta r \gg 1$）}
此时推迟效应不可忽略，场具有行波性质：
\begin{formulaBox}
\[
\begin{cases}
\dot{E}_\theta = j\sqrt{\dfrac{\mu_0}{\varepsilon_0}} \dfrac{\dot{I} \Delta l \sin\theta}{2\lambda r} e^{-j\beta r} \\[6pt]
\dot{H}_\varphi = j\dfrac{\dot{I} \Delta l \sin\theta}{2\lambda r} e^{-j\beta r} \\[6pt]
\dot{E}_r = \dot{E}_\varphi = \dot{H}_r = \dot{H}_\theta = 0
\end{cases}
\]
\end{formulaBox}

远区场的特点：
- 场量的相位随 $r$ 的增大不断滞后，具有行波性质；
- 等相位面为球面，故远区场是\textbf{球面波}；
- 电场和磁场同相位，存在向外传播的有功功率，形成电磁辐射。

\paragraph{远区场的关键参数}
\begin{itemize}
    \item \textbf{波阻抗}：自由空间中，远区电场与磁场的比值为常数：
    \[
    Z_0 = \frac{\dot{E}_\theta}{\dot{H}_\varphi} = \sqrt{\frac{\mu_0}{\varepsilon_0}} \approx 377\,\Omega
    \]
    波阻抗描述了远区场中电场和磁场的振幅关系。

    \item \textbf{辐射功率与辐射电阻}：穿过半径为 $r$ 的球面向外辐射的总电磁功率为：
    \[
    P = \oint_{A} \text{Re}(\dot{\vec{S}}) \cdot d\vec{A} = I^2 \left[ 80\pi^2 \left( \frac{\Delta l}{\lambda} \right)^2 \right] = I^2 R_e
    \]
    其中 $R_e = 80\pi^2 \left( \frac{\Delta l}{\lambda} \right)^2$ 称为\textbf{辐射电阻}，它表征了天线辐射电磁能量的能力。
\end{itemize}

\vspace{0.8cm}\hrule\vspace{0.8cm}
% 第6章 准静态电磁场
\section{准静态电磁场}

\subsection{电准静态场（EQS）}
当感应电场远小于库仑电场时，即 $\frac{\partial \vec{B}}{\partial t} \approx 0$，麦克斯韦方程组可简化为：
\begin{formulaBox}
\[
\begin{cases}
\nabla \times \vec{E} \approx 0 \\
\nabla \cdot \vec{D} = \rho_f \\
\nabla \times \vec{H} = \vec{J}_c + \frac{\partial \vec{D}}{\partial t} \\
\nabla \cdot \vec{B} = 0
\end{cases}
\]
\end{formulaBox}

与静态场相比，电准静态场具有与静电场类同的\textbf{有源无旋}特性，因此两种场的计算方法相同。

典型应用场景：
\begin{itemize}
    \item 电力传输系统和装置中的高压电场；
    \item 常用电子器件、设备附近的电场；
    \item 低频交流情况下，平板电容器中的电场。
\end{itemize}

\vspace{0.8cm}\hrule\vspace{0.8cm}
\subsection{磁准静态场（MQS）}
当位移电流远小于传导电流时，即 $\frac{\partial \vec{D}}{\partial t} \approx 0$，麦克斯韦方程组可简化为：
\begin{formulaBox}
\[
\begin{cases}
\nabla \times \vec{H} \approx \vec{J}_c \\
\nabla \cdot \vec{B} = 0 \\
\nabla \times \vec{E} = -\frac{\partial \vec{B}}{\partial t} \\
\nabla \cdot \vec{D} = \rho_f
\end{cases}
\]
\end{formulaBox}

与静态场相比，磁准静态场具有与静磁场类同的\textbf{有旋无源}特性，因此两种场的计算方法相同。

典型应用场景：
\begin{itemize}
    \item 低频电磁装置中的磁场问题（电机、变压器、感应加热、磁悬浮等）；
    \item 电工技术中的涡流问题；
    \item 电气工程中的大部分问题都属于磁准静态场问题。
\end{itemize}

磁准静态场中的位函数引出与一般时变电磁场相同：
\[
\begin{cases}
\nabla \cdot \vec{B} = 0 \implies \vec{B} = \nabla \times \vec{A} \\
\nabla \times \vec{E} = -\frac{\partial \vec{B}}{\partial t} \implies \vec{E} = -\frac{\partial \vec{A}}{\partial t} - \nabla \varphi
\end{cases}
\]
在库仑规范 $\nabla \cdot \vec{A} = 0$ 下，位函数满足：
\begin{formulaBox}
\[
\begin{cases}
\nabla^2 \vec{A} = -\mu \vec{J}_c \\
\nabla^2 \varphi = -\frac{\rho_f}{\varepsilon}
\end{cases}
\]
\end{formulaBox}

磁准静态场也称作\textbf{似稳场}，其核心近似是略去电磁场的波动性，认为场和源之间具有类似于静态场的“瞬时”对应关系。

\vspace{0.8cm}\hrule\vspace{0.8cm}
\subsection{似稳条件}
准静态近似的适用条件可分为两类：

\paragraph{1. 导体中的时变电磁场}
当满足 $\omega\varepsilon \ll \gamma$ 时，导体中的时变电磁场可按磁准静态场处理，这类场也称为\textbf{涡流场}。

\paragraph{2. 理想介质中的时变电磁场}
当满足 $\beta r \ll 1$（或 $r \ll \lambda$）时，理想介质中的时变电磁场可按磁准静态场处理。
物理含义：载流系统的尺寸必须远小于电磁波的波长，系统中各点的场可近似认为是同相位变化的。

\vspace{0.8cm}\hrule\vspace{0.8cm}
\subsection{集肤效应}
当交变电流流过导体时，电流密度分布呈现以下特点：
\begin{itemize}
    \item 靠近导体表面处电流密度大，深入导体内部电流密度小；
    \item 频率很高时，电流几乎只在导体表面附近的薄层中存在；
    \item 这种场量主要集中在导体表面附近的现象，称为\textbf{集肤效应}。
\end{itemize}

工程上常用\textbf{透入深度} $\delta$ 表示场量在良导体中的集肤程度，它等于场量振幅衰减到其表面值的 $1/e$ 时所经过的距离：
\begin{formulaBox}
\[
\delta = \frac{1}{\alpha} = \sqrt{\frac{2}{\omega\mu\gamma}}
\]
\end{formulaBox}

结论：\textbf{频率越高、导电性能越好的导体，集肤效应越显著}。

\vspace{0.8cm}\hrule\vspace{0.8cm}
\subsection{交流内阻抗}
在交流情况下，由于集肤效应，电流和电磁场集中在导体表面附近，导致有效载流截面积减小，因此导体的电阻和电感与直流时不同。

设导体中通过总电流 $\dot{I}$，其等效交流阻抗为 $Z = R + jX$，可通过坡印廷矢量的复数形式计算：
\begin{formulaBox}
\[
Z = \frac{-\oint_{A} (\dot{\vec{E}} \times \dot{\vec{H}}^*) \cdot d\vec{A}}{|\dot{I}|^2}
\]
其中：
\[
\begin{cases}
R = \frac{1}{|\dot{I}|^2} \text{Re}\left[ -\oint_{A} (\dot{\vec{E}} \times \dot{\vec{H}}^*) \cdot d\vec{A} \right] \\
X = \frac{1}{|\dot{I}|^2} \text{Im}\left[ -\oint_{A} (\dot{\vec{E}} \times \dot{\vec{H}}^*) \cdot d\vec{A} \right]
\end{cases}
\]
\end{formulaBox}

$R$ 为\textbf{交流电阻}（反映焦耳损耗），$X$ 为\textbf{内电抗}（反映导体内部磁场储能）。
\end{document}